% =====================================================================
%  CAPITULO 1 - continuacao
%  1.2 Lei de Coulomb  |  1.3 Campo eletrico
% =====================================================================

% =====================================================================
\subsection{Lei de Coulomb}
% =====================================================================

\begin{conceito}[Antes de começar]
\textbf{Lei de Coulomb} --- o módulo da força entre duas cargas puntiformes é
\[
F = \ke\,\frac{|q_1||q_2|}{d^2},\qquad
\ke = \SI{9e9}{\newton\meter\squared\per\coulomb\squared}\ \text{(no vácuo)}.
\]
A força é \emph{atrativa} para cargas de sinais opostos e \emph{repulsiva} para
sinais iguais, sempre na direção da reta que une as cargas.\par
\textbf{Princípio da superposição:} a força resultante sobre uma carga é a
\emph{soma vetorial} das forças que cada uma das demais exerce sobre ela.\par
\textbf{Dependências úteis:} $F \propto \dfrac{1}{d^2}$ (dobrar a distância divide
a força por 4) e $F \propto q_1 q_2$.
\end{conceito}

\legendaniveis

\blocofixacao

\begin{questao}[1]{}
Duas cargas $q_1 = \SI{+2}{\micro\coulomb}$ e $q_2 = \SI{-3}{\micro\coulomb}$
estão separadas por $\SI{4}{\centi\meter}$ no vácuo. Determine o módulo e o tipo
(atração ou repulsão) da força entre elas.
\dados{$\ke = \SI{9e9}{\newton\meter\squared\per\coulomb\squared}$.}
\resposta{$F = \SI{33.75}{\newton}$, atrativa}
\resolucao{Com $d = \SI{4}{\centi\meter} = \SI{0.04}{\meter}$:
\[
F=\ke\frac{|q_1||q_2|}{d^2}
=9\cdot10^9\cdot\frac{(\pot{2}{-6})(\pot{3}{-6})}{(0{,}04)^2}
=\SI{33.75}{\newton}.
\]
Como as cargas têm sinais opostos, a força é \textbf{atrativa}.}
\end{questao}

\begin{questao}[1]{}
Duas cargas $q_1 = q_2 = \SI{+5}{\micro\coulomb}$ estão a
$\SI{2}{\centi\meter}$ no vácuo. Determine o módulo e o tipo da força.
\dados{$\ke = \SI{9e9}{\newton\meter\squared\per\coulomb\squared}$.}
\resposta{$F = \SI{562.5}{\newton}$, repulsiva}
\resolucao{$F = 9\cdot10^9\cdot\dfrac{(\pot{5}{-6})^2}{(0{,}02)^2}
=\SI{562.5}{\newton}$, \textbf{repulsiva} (sinais iguais).}
\end{questao}

\begin{questao}[1]{}
Duas cargas puntiformes se repelem com força de intensidade
$\pot{42}{-5}\,\si{\newton}$. Reduzindo-se a distância entre elas à
\textbf{metade}, a nova força de repulsão será:
\begin{alternativas}[2]
\item $\pot{10.5}{-5}\,\si{\newton}$
\item $\pot{21}{-5}\,\si{\newton}$
\item $\pot{84}{-5}\,\si{\newton}$
\item $\pot{168}{-5}\,\si{\newton}$
\item inalterada
\end{alternativas}
\resposta{(d)}
\resolucao{Como $F \propto \dfrac{1}{d^2}$, reduzir $d$ à metade \emph{multiplica}
a força por 4:
\[
F' = 4F = 4\cdot\pot{42}{-5} = \pot{168}{-5}\,\si{\newton}.
\]}
\end{questao}

\blocoaplicacao

\begin{questao}[2]{}
A força elétrica entre duas partículas de cargas $q$ e $q/2$, separadas por uma
distância $d$ no vácuo, é $F$. A força entre duas partículas de cargas $q$ e $2q$,
separadas por $d/2$, também no vácuo, é:
\begin{alternativas}[3]
\item $F$
\item $2F$
\item $4F$
\item $8F$
\item $16F$
\end{alternativas}
\resposta{(e)}
\resolucao{Situação inicial: $F = \ke\dfrac{q\cdot(q/2)}{d^2} = \dfrac{\ke q^2}{2d^2}$.\\
Situação nova: $F' = \ke\dfrac{q\cdot 2q}{(d/2)^2} = \dfrac{2\ke q^2}{d^2/4}
= \dfrac{8\ke q^2}{d^2}$.\\
Dividindo: $\dfrac{F'}{F} = \dfrac{8\ke q^2/d^2}{\ke q^2/(2d^2)} = 16$, logo
$F' = 16F$. O produto das cargas quadruplicou e a distância caiu à metade
(fator 4 na força): $4\times4 = 16$.}
\end{questao}

\begin{questao}[2]{}
Se a força entre dois objetos carregados se mantém \emph{constante} mesmo quando a
carga de cada um é reduzida à metade, então a distância entre eles:
\begin{alternativas}
\item foi quadruplicada.
\item foi duplicada.
\item foi reduzida à quarta parte.
\item foi reduzida à metade.
\item permaneceu constante.
\end{alternativas}
\resposta{(d)}
\resolucao{Reduzir cada carga à metade divide o produto $q_1 q_2$ por 4, o que
sozinho reduziria $F$ a $1/4$. Para $F$ permanecer constante, o fator $1/d^2$
precisa \emph{quadruplicar}, ou seja, $d^2$ deve cair a $1/4$, o que significa $d$
reduzido à \textbf{metade}:
\[
F=\ke\frac{(q_1/2)(q_2/2)}{d'^2}=\ke\frac{q_1q_2}{4d'^2}
\stackrel{!}{=}\ke\frac{q_1q_2}{d^2}
\ \Rightarrow\ d'^2=\frac{d^2}{4}\ \Rightarrow\ d'=\frac{d}{2}.
\]}
\end{questao}

\begin{questao}[2]{}
Três cargas iguais de $\SI{+1}{\micro\coulomb}$ ocupam os vértices de um triângulo
equilátero de lado $\SI{5}{\centi\meter}$. Determine o módulo da força resultante
sobre uma delas.

\circuitoq{%
\begin{tikzpicture}[scale=1,line width=0.8pt,>=Stealth]
 \coordinate (A) at (0,0);
 \coordinate (B) at (3,0);
 \coordinate (C) at (1.5,2.6);
 \foreach \p/\n in {A/A,B/B,C/C}{
   \shade[ball color=Vcol!25] (\p) circle (0.28);
   \draw[Vcol] (\p) circle (0.28);
   \node at (\p) {\footnotesize$+q$};}
 \draw[dashed,cinza] (A)--(B)--(C)--cycle;
 \node[below] at (1.5,-0.35) {$\SI{5}{\centi\meter}$};
\end{tikzpicture}}{Três cargas iguais nos vértices de um triângulo equilátero.}

\resposta{$F_R \approx \SI{6.24}{\newton}$}
\resolucao{Cada par de cargas exerce uma força de módulo
\[
F=\ke\frac{q^2}{d^2}=9\cdot10^9\cdot\frac{(\pot{1}{-6})^2}{(0{,}05)^2}
=\SI{3.6}{\newton}.
\]
Sobre a carga escolhida atuam \emph{duas} forças iguais de \SI{3.6}{\newton}, com
$\SI{60}{\degree}$ entre elas (geometria do equilátero). A resultante é
\[
F_R = \sqrt{F^2+F^2+2F^2\cos 60^\circ}=F\sqrt{3}
=3{,}6\sqrt{3}\approx\SI{6.24}{\newton},
\]
dirigida para fora do triângulo, ao longo da bissetriz do vértice.}
\end{questao}

\begin{questao}[2]{}
Quatro cargas positivas iguais de $\SI{1}{\micro\coulomb}$ ocupam os vértices de
um quadrado de lado $\SI{10}{\centi\meter}$. Determine o módulo da força
resultante sobre uma delas.
\dados{$\ke = \SI{9e9}{\newton\meter\squared\per\coulomb\squared}$; $\sqrt2\approx1{,}41$.}

\circuitoq{%
\begin{tikzpicture}[scale=1,line width=0.8pt,>=Stealth]
 \foreach \p in {(0,0),(2.4,0),(2.4,2.4),(0,2.4)}{
   \shade[ball color=Vcol!25] \p circle (0.22);}
 \node at (0,0){\footnotesize$+q$};\node at (2.4,0){\footnotesize$+q$};
 \node at (2.4,2.4){\footnotesize$+q$};\node at (0,2.4){\footnotesize$+q$};
 \draw[dashed,cinza] (0,0)--(2.4,0)--(2.4,2.4)--(0,2.4)--cycle;
 \draw[dashed,cinza] (0,0)--(2.4,2.4);
 \node[below] at (1.2,-0.3){$\SI{10}{\centi\meter}$};
\end{tikzpicture}}{Quatro cargas iguais nos vértices de um quadrado.}

\resposta{$F_R \approx \SI{1.72}{\newton}$}
\resolucao{Sobre a carga de um vértice atuam três forças repulsivas:
duas dos vértices \emph{adjacentes} (ao longo dos lados, perpendiculares entre si)
e uma do vértice \emph{oposto} (ao longo da diagonal).\\
\textbf{Força de um lado} ($d = \SI{10}{\centi\meter}$):
$F_L = \ke\dfrac{q^2}{d^2} = 9\cdot10^9\dfrac{(\pot{1}{-6})^2}{(0{,}1)^2}
= \SI{0.9}{\newton}$.\\
\textbf{Força da diagonal} ($d' = d\sqrt2$):
$F_D = \ke\dfrac{q^2}{(d\sqrt2)^2} = \dfrac{F_L}{2} = \SI{0.45}{\newton}$.\\
As duas forças de lado, perpendiculares, têm resultante
$F_L\sqrt2 \approx \SI{1.27}{\newton}$, na direção da diagonal --- \emph{mesma}
direção de $F_D$. Somando:
\[
F_R = F_L\sqrt2 + F_D \approx 1{,}27 + 0{,}45 = \SI{1.72}{\newton}.
\]}
\end{questao}

\begin{questao}[2]{}
Uma pequena esfera de massa $\SI{2}{\gram}$, eletrizada com carga $q$, permanece
em equilíbrio, suspensa, sob a ação combinada do peso e de um campo elétrico
uniforme vertical de intensidade $\SI{e5}{\newton\per\coulomb}$, dirigido para
cima. Determine a carga $q$.
\dados{$g = \SI{10}{\meter\per\second\squared}$.}
\resposta{$q = \SI{0.2}{\micro\coulomb}$}
\resolucao{No equilíbrio, a força elétrica (para cima) equilibra o peso (para
baixo): $qE = mg$, logo
\[
q=\frac{mg}{E}=\frac{\pot{2}{-3}\cdot10}{\pot{1}{5}}
=\pot{2}{-7}\,\si{\coulomb}=\SI{0.2}{\micro\coulomb}.
\]
Para a força ser para cima (mesmo sentido de $E$), a carga deve ser
\emph{positiva}. É o princípio do experimento de Millikan, que mediu a carga do
elétron equilibrando gotas de óleo.}
\end{questao}

\blocoaprofundamento

\begin{questao}[3]{}
Duas cargas $Q_1$ e $Q_2 = 4Q_1$ (ambas positivas) estão fixas nos pontos A e B,
distantes $\SI{30}{\centi\meter}$. Em que posição $x$, medida a partir de A sobre
a reta AB, deve-se colocar uma carga $Q_3$ para que ela fique em equilíbrio sob
ação apenas de forças elétricas?

\circuitoq{%
\begin{tikzpicture}[scale=1,line width=0.8pt]
 \draw[cinza] (0,0)--(6,0);
 \shade[ball color=Vcol!25] (0,0) circle (0.25);
 \shade[ball color=Vcol!25] (6,0) circle (0.25);
 \node at (0,0) {\footnotesize$Q_1$}; \node at (6,0) {\footnotesize$Q_2$};
 \node[below] at (0,-0.35) {A}; \node[below] at (6,-0.35) {B};
 \fill[laranja] (2,0) circle (0.09);
 \node[above,laranja] at (2,0.15) {$Q_3$};
 \draw[<->,cinza] (0,-0.9)--(2,-0.9) node[midway,below]{$x$};
 \draw[<->,cinza] (0,-1.5)--(6,-1.5) node[midway,below]{$\SI{30}{\centi\meter}$};
\end{tikzpicture}}{Posição de equilíbrio entre duas cargas.}

\resposta{$x = \SI{10}{\centi\meter}$ a partir de A}
\resolucao{No equilíbrio, as forças que $Q_1$ e $Q_2$ exercem sobre $Q_3$ têm
mesmo módulo e sentidos opostos (o que só é possível \emph{entre} as cargas, já
que ambas são positivas). Igualando os módulos, com $d = \SI{30}{\centi\meter}$:
\[
\ke\frac{Q_1 Q_3}{x^2}=\ke\frac{4Q_1\,Q_3}{(d-x)^2}
\ \Rightarrow\
\frac{1}{x^2}=\frac{4}{(d-x)^2}
\ \Rightarrow\
\frac{d-x}{x}=2.
\]
Logo $d - x = 2x \Rightarrow x = \dfrac{d}{3} = \SI{10}{\centi\meter}$. O ponto de
equilíbrio fica \emph{mais próximo} da carga menor ($Q_1$), como esperado.
Note que o valor de $Q_3$ e seu sinal não influenciam a posição.}
\end{questao}

\begin{questao}[3]{}
Quatro cargas ocupam os vértices de um quadrado: duas delas são $+Q$ e $-Q$ (mesmo
módulo), e as outras duas, $q_1$ e $q_2$, são desconhecidas. Coloca-se uma carga
de prova positiva no centro e observa-se que a força resultante sobre ela aponta
na direção da diagonal, \emph{do vértice $-Q$ para o vértice $+Q$}. Sobre $q_1$ e
$q_2$ (nos outros dois vértices), pode-se afirmar que, para que essa força
resultante fique exatamente sobre a diagonal $-Q\,$--$\,+Q$:
\begin{alternativas}
\item devem ser iguais em módulo e de mesmo sinal entre si.
\item devem ser iguais em módulo e de sinais opostos.
\item devem ser ambas nulas.
\item devem ser ambas positivas e diferentes.
\item nada se pode afirmar.
\end{alternativas}
\resposta{(a)}
\resolucao{A força devida ao par $+Q$/$-Q$ já está sobre a diagonal
$-Q\,$--$\,+Q$. Para que a resultante \emph{total} continue sobre essa mesma
diagonal, a contribuição do par $q_1$/$q_2$ (na outra diagonal) não pode ter
componente perpendicular a ela: as forças de $q_1$ e $q_2$ sobre a carga central
devem se \emph{cancelar} nessa direção transversal. Isso ocorre quando $q_1$ e
$q_2$ têm o mesmo módulo e o mesmo sinal (por simetria em relação à diagonal
principal), produzindo forças opostas sobre a carga de prova que se anulam.}
\end{questao}

\begin{questao}[3]{}
Duas cargas fixas, $q_1 = \SI{+9}{\micro\coulomb}$ em $x=0$ e
$q_2 = \SI{+4}{\micro\coulomb}$ em $x = \SI{10}{\centi\meter}$, estão sobre o eixo
$x$. Uma terceira carga $q_3$ é colocada sobre esse eixo.
\begin{itensq}
\item Determine a posição em que $q_3$ fica em equilíbrio.
\item Mostre que existe uma \emph{segunda} raiz matemática e explique por que ela
      deve ser descartada.
\item O equilíbrio encontrado é \emph{estável} ou \emph{instável}? Analise para
      $q_3$ positiva e negativa.
\end{itensq}
\resposta{(a) $x = \SI{6}{\centi\meter}$; (b) $x=\SI{30}{\centi\meter}$ (fora do
segmento); (c) estável para $q_3<0$, instável para $q_3>0$}
\resolucao{(a) Igualando os módulos das forças (a carga $q_3$ cancela-se):
\[
\frac{\ke\,9}{x^{2}}=\frac{\ke\,4}{(10-x)^{2}}
\;\Longrightarrow\;
\frac{3}{x}=\frac{2}{10-x}
\;\Longrightarrow\;
30-3x=2x
\;\Longrightarrow\;
x=\SI{6}{\centi\meter}.
\]
(b) Ao elevar ao quadrado, a equação admite também $\dfrac{3}{x}=-\dfrac{2}{10-x}$,
que fornece $x = \SI{30}{\centi\meter}$. Nesse ponto, \emph{fora} do segmento,
as duas forças apontam no \textbf{mesmo sentido} (ambas as cargas são positivas e
empurram $q_3$ para longe) --- elas se somam e nunca se cancelam. A raiz é
espúria, introduzida pela elevação ao quadrado.

(c) \textbf{Análise de estabilidade.} Suponha $q_3$ deslocada ligeiramente para a
direita (aproximando-se de $q_2$):
\begin{itemize}[leftmargin=1.6em,itemsep=1pt,topsep=2pt]
\item Se $q_3 > 0$: a repulsão de $q_2$ (agora mais próxima) supera a de $q_1$, e
      $q_3$ é empurrada \emph{ainda mais} para longe do ponto de equilíbrio ---
      \textbf{instável}.
\item Se $q_3 < 0$: a atração por $q_2$ cresce, puxando $q_3$ \emph{de volta}?
      Não --- a atração por $q_2$ a puxa para longe de $x=6$. Nesse eixo,
      portanto, também é instável.
\end{itemize}
\textbf{Análise completa:} para $q_3<0$, o equilíbrio é estável na direção
\emph{transversal} ao eixo (qualquer deslocamento lateral gera força restauradora),
mas instável ao longo do eixo. Para $q_3>0$, ocorre o oposto.\\
\textbf{Conclusão profunda --- teorema de Earnshaw:} \emph{nenhuma} configuração
de cargas estáticas produz equilíbrio estável em \textbf{todas} as direções. Sempre
há ao menos uma direção instável. É por isso que não se consegue levitar um objeto
apenas com cargas fixas --- são necessários campos variáveis (armadilhas de Paul)
ou diamagnetismo.}
\end{questao}

\begin{questao}[3]{}
Compare a força elétrica e a força gravitacional entre dois elétrons.
\begin{itensq}
\item Mostre que a razão $F_e/F_g$ \emph{independe} da distância.
\item Calcule seu valor numérico.
\item Que conclusão isso permite sobre a estrutura da matéria e sobre a
      astronomia?
\end{itensq}
\dados{$\ke=\SI{9e9}{}$; $G=\SI{6.67e-11}{}$; $e=\SI{1.6e-19}{\coulomb}$;
$m_e=\SI{9.11e-31}{\kilogram}$ (SI).}
\resposta{(a) demonstração; (b) $\approx\pot{4.2}{42}$; (c) ver comentário}
\resolucao{(a) Ambas as forças seguem a lei do inverso do quadrado:
\[
F_e=\ke\frac{e^{2}}{d^{2}},
\qquad
F_g=G\frac{m_e^{2}}{d^{2}}
\;\Longrightarrow\;
\frac{F_e}{F_g}=\frac{\ke\,e^{2}}{G\,m_e^{2}}.
\]
O fator $d^{2}$ \textbf{cancela-se}: a razão é uma constante universal.\\
(b) Numericamente:
\[
\frac{F_e}{F_g}=\frac{(9\cdot10^{9})(\pot{1.6}{-19})^{2}}
{(\pot{6.67}{-11})(\pot{9.11}{-31})^{2}}\approx\pot{4.2}{42}.
\]
(c) \textbf{Duas conclusões profundas:}
\begin{itemize}[leftmargin=1.6em,itemsep=2pt,topsep=3pt]
\item \textbf{Na escala atômica, a gravidade é irrelevante.} A força elétrica é
      $10^{42}$ vezes maior --- por isso a estrutura de átomos, moléculas,
      cristais e de toda a química é governada \emph{exclusivamente} por forças
      eletromagnéticas. Nenhum modelo atômico precisa incluir gravidade.
\item \textbf{Na escala astronômica, ocorre o oposto.} Corpos celestes são
      eletricamente \emph{neutros} (as cargas positivas e negativas se cancelam
      quase perfeitamente), então a força elétrica resultante é praticamente
      nula. Já a massa só se \emph{acumula} --- não existe "massa negativa". Por
      isso a gravidade, apesar de absurdamente mais fraca, domina o cosmos.
\end{itemize}
\emph{A natureza usa a força fraca para construir galáxias e a força forte para
construir átomos --- precisamente porque uma se cancela e a outra não.}}
\end{questao}

\begin{questao}[3]{}
Duas pequenas esferas idênticas, de massa $\SI{0.5}{\gram}$ cada, são suspensas
do mesmo ponto por fios isolantes de $\SI{20}{\centi\meter}$. Ao receberem cargas
iguais $q$, elas se afastam até que cada fio forme $\SI{10}{\degree}$ com a
vertical.

\circuitoq{%
\begin{tikzpicture}[scale=1,line width=0.8pt,>=Stealth]
 \draw[cinza!60,line width=2pt] (-1.6,3)--(1.6,3);
 \coordinate (O) at (0,3);
 \coordinate (P1) at ({-2.4*sin(10)},{3-2.4*cos(10)});
 \coordinate (P2) at ({ 2.4*sin(10)},{3-2.4*cos(10)});
 \draw[cinza,thick] (O)--(P1); \draw[cinza,thick] (O)--(P2);
 \draw[dashed,cinza!70] (O)--(0,0.3);
 \shade[ball color=Vcol!30] (P1) circle (0.2);
 \shade[ball color=Vcol!30] (P2) circle (0.2);
 \node[font=\footnotesize] at (0.34,2.55) {$10^\circ$};
 \node[left,font=\footnotesize] at ($(P1)+(-0.25,0)$) {$q$};
 \node[right,font=\footnotesize] at ($(P2)+(0.25,0)$) {$q$};
 \draw[<->,cinza] ($(P1)+(0,-0.45)$)--($(P2)+(0,-0.45)$)
       node[midway,below,font=\footnotesize]{$d$};
 \node[right,font=\footnotesize,cinza] at (0.9,1.9) {$L=\SI{20}{\centi\meter}$};
\end{tikzpicture}}{Pêndulos eletrostáticos idênticos.}

\begin{itensq}
\item Determine a separação $d$ entre as esferas.
\item Calcule a força elétrica entre elas.
\item Determine a carga $q$ de cada esfera.
\end{itensq}
\dados{$g=\SI{10}{\meter\per\second\squared}$; $\sin10^\circ=0{,}174$;
$\tan10^\circ=0{,}176$; $\ke=\SI{9e9}{}$.}
\resposta{(a) $d\approx\SI{6.9}{\centi\meter}$; (b) $F\approx\SI{0.88}{\milli\newton}$;
(c) $q\approx\SI{22}{\nano\coulomb}$}
\resolucao{(a) Cada esfera se afasta $L\sin\theta$ da vertical, logo
\[
d=2L\sin\theta=2(0{,}20)(0{,}174)\approx\SI{0.069}{\meter}=\SI{6.9}{\centi\meter}.
\]
(b) No equilíbrio de cada esfera (tração, peso e força elétrica),
a decomposição fornece $\tan\theta = \dfrac{F_e}{mg}$:
\[
F_e = mg\tan\theta=(\pot{0.5}{-3})(10)(0{,}176)
\approx\pot{8.8}{-4}\,\si{\newton}=\SI{0.88}{\milli\newton}.
\]
(c) Da lei de Coulomb, com as duas cargas iguais:
\[
F_e=\ke\frac{q^{2}}{d^{2}}
\;\Longrightarrow\;
q=d\sqrt{\frac{F_e}{\ke}}
=0{,}069\sqrt{\frac{\pot{8.8}{-4}}{9\cdot10^{9}}}
\approx\pot{2.2}{-8}\,\si{\coulomb}=\SI{22}{\nano\coulomb}.
\]
\textbf{Ordem de grandeza:} apenas $\sim\pot{1.4}{11}$ elétrons em excesso ---
uma quantidade ínfima --- já produz um afastamento visível. \emph{Isso mostra
quão intensa é a força elétrica em escala macroscópica.}}
\end{questao}

\begin{questao}[3]{}
Quatro cargas de mesmo módulo $q=\SI{1}{\micro\coulomb}$ ocupam os vértices de um
quadrado de lado $a=\SI{10}{\centi\meter}$, \textbf{alternando os sinais}
($+,-,+,-$ ao percorrer o perímetro). Determine o módulo e a direção da força
resultante sobre uma das cargas positivas.
\dados{$\ke=\SI{9e9}{}$; $\sqrt2\approx1{,}41$.}
\resposta{$F_R \approx \SI{0.82}{\newton}$, dirigida ao centro do quadrado}
\resolucao{\textbf{Forças envolvidas.} Sobre a carga positiva escolhida atuam:
\begin{itemize}[leftmargin=1.6em,itemsep=1pt,topsep=2pt]
\item duas forças dos vizinhos \emph{adjacentes} (que são \textbf{negativos}):
      são de \textbf{atração}, ao longo dos lados;
\item uma força da carga da \emph{diagonal} (que é \textbf{positiva}):
      \textbf{repulsão}, ao longo da diagonal.
\end{itemize}
\textbf{Módulos:}
\[
F_L=\ke\frac{q^{2}}{a^{2}}=9\cdot10^{9}\frac{(\pot{1}{-6})^{2}}{(0{,}1)^{2}}
=\SI{0.9}{\newton},
\qquad
F_D=\ke\frac{q^{2}}{(a\sqrt2)^{2}}=\frac{F_L}{2}=\SI{0.45}{\newton}.
\]
\textbf{Composição.} As duas atrações, perpendiculares entre si, têm resultante
$F_L\sqrt2 \approx \SI{1.27}{\newton}$ apontando \emph{para o centro} do quadrado
(ao longo da diagonal). A repulsão da diagonal aponta \emph{para fora}, na mesma
reta. Logo elas se \textbf{subtraem}:
\[
F_R=F_L\sqrt2-F_D = 1{,}27-0{,}45\approx\SI{0.82}{\newton},
\]
dirigida \textbf{para o centro} do quadrado.\\
\textbf{Contraste com o caso de cargas iguais} (visto anteriormente): lá as três
forças eram repulsivas e se \emph{somavam} ($F_R=F_L\sqrt2+F_D=\SI{1.72}{\newton}$),
apontando para fora. Aqui, alternar os sinais inverte o sentido e reduz o módulo
para menos da metade. \emph{A geometria é a mesma; a configuração de sinais muda
tudo.}}
\end{questao}

\begin{questao}[3]{}
Duas cargas fixas, $q_1=+q$ e $q_2=+4q$, estão separadas por $d$. Deseja-se
colocar uma terceira carga $q_3$ de modo que \textbf{as três} fiquem em
equilíbrio simultaneamente.
\begin{itensq}
\item Determine a posição de $q_3$.
\item Determine o valor e o sinal de $q_3$.
\item Esse equilíbrio é estável?
\end{itensq}
\resposta{(a) a $d/3$ de $q_1$; (b) $q_3=-\dfrac{4q}{9}$; (c) instável}
\resolucao{(a) Para $q_3$ estar em equilíbrio, ela deve ficar onde os campos de
$q_1$ e $q_2$ se cancelam (calculado anteriormente):
\[
\frac{q}{x^{2}}=\frac{4q}{(d-x)^{2}}
\;\Rightarrow\;
\frac{d-x}{x}=2
\;\Rightarrow\;
x=\frac{d}{3}.
\]
(b) Mas isso garante apenas o equilíbrio de $q_3$. Para que \emph{$q_1$} também
esteja equilibrada, a repulsão de $q_2$ sobre ela deve ser cancelada pela atração
de $q_3$:
\[
\underbrace{\ke\frac{q\cdot 4q}{d^{2}}}_{\text{repulsão de }q_2}
=\underbrace{\ke\frac{q\,|q_3|}{(d/3)^{2}}}_{\text{atração de }q_3}
\;\Longrightarrow\;
\frac{4q}{d^{2}}=\frac{9|q_3|}{d^{2}}
\;\Longrightarrow\;
|q_3|=\frac{4q}{9}.
\]
O sinal deve ser \textbf{negativo} (para atrair), logo
$q_3 = -\dfrac{4q}{9}$.\\
\textbf{Verificação em $q_2$:} a repulsão de $q_1$ sobre $q_2$ é
$\ke\dfrac{4q^2}{d^2}$; a atração de $q_3$ (à distância $2d/3$) é
$\ke\dfrac{4q\cdot(4q/9)}{(2d/3)^{2}}=\ke\dfrac{(16q^2/9)}{(4d^2/9)}
=\ke\dfrac{4q^{2}}{d^{2}}$. \checkmark\ São iguais --- as três estão em
equilíbrio.\\
(c) \textbf{Instável}, pelo teorema de Earnshaw: qualquer deslocamento infinitesimal
de qualquer das cargas destrói o equilíbrio, e as forças resultantes \emph{afastam}
ainda mais o sistema da configuração original. Não existe arranjo eletrostático
puro que seja estável.}
\end{questao}

\begin{questao}[3]{}
Uma pequena esfera de massa $m = \SI{10}{\gram}$, eletrizada com carga $q$, é
suspensa por um fio isolante de massa desprezível em uma região de campo elétrico
\emph{horizontal} e uniforme, de intensidade $E = \SI{e4}{\newton\per\coulomb}$.
No equilíbrio, o fio forma $\SI{30}{\degree}$ com a vertical.

\circuitoq{%
\begin{tikzpicture}[scale=1,line width=0.8pt,>=Stealth]
 \draw[cinza!60,line width=2pt] (-2,2.6)--(2,2.6);
 \coordinate (O) at (0,2.6);
 \coordinate (P) at ({2.2*sin(30)},{2.6-2.2*cos(30)});
 \draw[cinza,thick] (O)--(P);
 \draw[dashed,cinza!70] (O)--(0,0.2);
 \shade[ball color=Vcol!30] (P) circle (0.22);
 \node[right,font=\footnotesize] at ($(P)+(0.28,0)$) {$m,\,q$};
 \draw[->,Vcol,thick] (P)--($(P)+(1.0,0)$) node[right,font=\footnotesize]{$\vec F_e$};
 \draw[->,cinza,thick] (P)--($(P)+(0,-0.9)$) node[below,font=\footnotesize]{$\vec P$};
 \node[font=\footnotesize] at (0.42,1.9) {$30^\circ$};
 \foreach \y in {0.4,1.0,1.6} \draw[->,Vcol!45] (-1.9,\y)--(-1.2,\y);
 \node[Vcol,font=\footnotesize] at (-1.55,2.05) {$\vec E$};
\end{tikzpicture}}{Pêndulo eletrostático em campo horizontal.}

\begin{itensq}
\item Determine a carga $q$ da esfera.
\item Calcule a tração no fio.
\item Se o campo dobrasse, o ângulo dobraria? Justifique quantitativamente.
\end{itensq}
\dados{$g = \SI{10}{\meter\per\second\squared}$; $\tan30^\circ\approx0{,}577$.}
\resposta{(a) $q \approx \SI{5.8}{\micro\coulomb}$; (b) $T\approx\SI{0.115}{\newton}$;
(c) não --- o ângulo cresce menos que proporcionalmente}
\resolucao{\textbf{(a)} No equilíbrio, decompondo a tração:
\[
T\sin\theta = F_e = qE,
\qquad
T\cos\theta = P = mg
\quad\Longrightarrow\quad
\tan\theta=\frac{qE}{mg}.
\]
Isolando $q$:
\[
q=\frac{mg\tan\theta}{E}
=\frac{(\pot{10}{-3})(10)(0{,}577)}{\pot{1}{4}}
\approx\pot{5.8}{-6}\,\si{\coulomb}=\SI{5.8}{\micro\coulomb}.
\]
\textbf{(b)} $T = \dfrac{mg}{\cos\theta} = \dfrac{0{,}1}{0{,}866}
\approx\SI{0.115}{\newton}$.\\
\textbf{(c)} \textbf{Não.} A relação é $\tan\theta \propto E$, e \emph{não}
$\theta \propto E$. Dobrando o campo:
\[
\tan\theta' = 2\tan30^\circ = 1{,}155
\;\Rightarrow\;
\theta' = \arctan(1{,}155)\approx 49{,}1^\circ,
\]
que é \emph{menos} que $2\times30^\circ = 60^\circ$. A função tangente cresce mais
rápido que o ângulo, então o ângulo responde de forma \emph{sublinear}. No limite
$E\to\infty$, $\theta\to90^\circ$ mas nunca o atinge --- o fio jamais fica
horizontal.}
\end{questao}

\begin{questao}[3]{}
Três cargas puntiformes estão sobre uma reta: $+Q$ na origem, $-2Q$ em $x=a$ e
$+Q$ em $x=2a$ (um \emph{quadrupolo linear}).
\begin{itensq}
\item Determine a força resultante sobre a carga central.
\item Determine a força resultante sobre a carga da extremidade em $x=0$.
\item Explique por que o sistema, embora tenha carga total nula, ainda exerce
      força sobre uma carga externa distante.
\end{itensq}
\resposta{(a) nula, por simetria; (b) $F=\dfrac{7\ke Q^{2}}{4a^{2}}$ dirigida para
$+x$; (c) ver comentário}
\resolucao{(a) A carga central $-2Q$ está \emph{equidistante} das duas cargas
$+Q$, que são iguais. As duas atrações têm mesmo módulo e sentidos opostos: a
resultante é \textbf{nula} por simetria.\\
(b) Sobre a carga em $x=0$ atuam:
\begin{itemize}[leftmargin=1.6em,itemsep=1pt,topsep=2pt]
\item atração de $-2Q$ (distância $a$):
      $F_1=\ke\dfrac{Q\cdot2Q}{a^{2}}=\dfrac{2\ke Q^{2}}{a^{2}}$, para $+x$;
\item repulsão de $+Q$ (distância $2a$):
      $F_2=\ke\dfrac{Q^{2}}{4a^{2}}$, para $-x$.
\end{itemize}
Resultante:
\[
F=F_1-F_2=\frac{2\ke Q^{2}}{a^{2}}-\frac{\ke Q^{2}}{4a^{2}}
=\frac{\ke Q^{2}}{a^{2}}\left(2-\frac14\right)
=\frac{7\ke Q^{2}}{4a^{2}},
\]
dirigida para $+x$ (para dentro do conjunto).\\
(c) \textbf{Por que a carga total nula não implica campo nulo.} A carga total é
$Q-2Q+Q=0$, então \emph{a longas distâncias} o campo decai muito mais rápido que
o de uma carga isolada. Mas ele não é nulo: as cargas ocupam \emph{posições
diferentes}, e essa separação espacial cria momentos de ordem superior.\\
\textbf{Hierarquia dos decaimentos:}
\begin{center}\small\sffamily
\begin{tabular}{@{}l c c@{}}
\toprule
\textbf{Configuração} & \textbf{Carga total} & \textbf{Campo distante}\\
\midrule
Monopolo (carga isolada) & $\neq0$ & $\propto 1/r^{2}$\\
Dipolo ($+q,-q$)         & $0$     & $\propto 1/r^{3}$\\
Quadrupolo (este caso)   & $0$     & $\propto 1/r^{4}$\\
\bottomrule
\end{tabular}
\end{center}
\emph{Cada cancelamento sucessivo acelera o decaimento --- é o princípio da
expansão multipolar, base da análise de antenas e de moléculas polares.}}
\end{questao}

\blocodesafio

\begin{questao}[4]{}
Uma carga $Q$ é dividida em duas partes, $q$ e $(Q-q)$, que são então colocadas a
uma distância fixa $d$ uma da outra.
\begin{itensq}
\item Obtenha a expressão da força entre elas em função de $q$.
\item Determine o valor de $q$ que \textbf{maximiza} essa força.
\item Qual é a força máxima?
\end{itensq}
\resposta{(a) $F=\dfrac{\ke\,q(Q-q)}{d^{2}}$; (b) $q=Q/2$;
(c) $F_{\max}=\dfrac{\ke Q^{2}}{4d^{2}}$}
\resolucao{(a) Pela lei de Coulomb,
\[
F(q)=\ke\,\frac{q\,(Q-q)}{d^{2}}=\frac{\ke}{d^{2}}\left(Qq-q^{2}\right).
\]
(b) É uma \textbf{parábola} com concavidade para baixo em função de $q$. O máximo
ocorre no vértice. Derivando:
\[
\frac{\mathrm{d}F}{\mathrm{d}q}=\frac{\ke}{d^{2}}\left(Q-2q\right)=0
\;\Longrightarrow\;
\boxed{\,q=\frac{Q}{2}\,}
\]
(Sem cálculo: o vértice de $Qq-q^2$ está em $q = -\dfrac{b}{2a}=\dfrac{Q}{2}$.)\\
(c) Substituindo:
\[
F_{\max}=\frac{\ke}{d^{2}}\cdot\frac{Q}{2}\cdot\frac{Q}{2}
=\frac{\ke\,Q^{2}}{4d^{2}}.
\]
\textbf{Interpretação.} A força é máxima quando a carga é dividida
\emph{igualmente}. Isso decorre da desigualdade entre médias: para soma fixa
($q+(Q-q)=Q$), o produto é máximo quando as parcelas são iguais.\\
\textbf{Casos-limite:} se $q\to0$ ou $q\to Q$, uma das partes fica sem carga e
$F\to0$ --- coerente. \emph{Este é um problema de otimização disfarçado de
eletrostática: reconhecer a estrutura matemática vale mais que decorar a
fórmula.}}
\end{questao}

% BEGIN BANCO COMPLEMENTAR Lei de Coulomb
% =====================================================================
%  Banco complementar --- Lei de Coulomb  (REVISADO)
%  Substitui a versao gerada por template: 32 questoes, todas com
%  enunciado distinto. Nivel 4 reclassificado conforme o criterio da
%  obra (sintese, demonstracao, otimizacao ou modelagem nao rotineira).
%  Distribuicao mantida: 7 N1 + 10 N2 + 6 N3 + 9 N4 = 32
% =====================================================================

\subsubsection*{Banco complementar --- Lei de Coulomb}

\begin{atencao}[Organização do banco]
As questões abaixo complementam o tópico sem repetir enunciados da seção
principal. Cada nível explora uma \emph{estrutura de raciocínio} diferente:
N1 aplica a fórmula em suas quatro formas isoladas ($F$, $q$, $r$, meio);
N2 exige composição ou reinterpretação; N3 combina vetores, equilíbrio e
redistribuição de carga; N4 pede demonstração, integração, otimização ou
modelagem.
\end{atencao}

\blocofixacao

\begin{questao}[1]{}
Duas cargas puntiformes $q_1=\SI{20}{\nano\coulomb}$ e
$q_2=\SI{30}{\nano\coulomb}$ estão separadas por $\SI{3}{\centi\meter}$ no
vácuo. Determine o módulo da força elétrica e classifique-a.
\dados{$\ke=\SI{9e9}{\newton\meter\squared\per\coulomb\squared}$}
\resposta{$F=\SI{6}{\milli\newton}$, repulsiva}
\resolucao{Convertendo antes de substituir --- este é o erro mais comum:
$q_1=\pot{2}{-8}\,\si{\coulomb}$, $q_2=\pot{3}{-8}\,\si{\coulomb}$,
$r=\pot{3}{-2}\,\si{\meter}$.
\[
F=\ke\frac{q_1q_2}{r^{2}}
 =\frac{\pot{9}{9}\cdot\pot{2}{-8}\cdot\pot{3}{-8}}{(\pot{3}{-2})^{2}}
 =\frac{\pot{5.4}{-6}}{\pot{9}{-4}}=\pot{6}{-3}\,\si{\newton}.
\]
Ambas positivas $\Rightarrow$ repulsão.}
\end{questao}

\begin{questao}[1]{}
Duas esferas idênticas recebem cargas iguais $q$ e se repelem com força de
$\SI{0.90}{\newton}$ quando separadas por $\SI{20}{\centi\meter}$.
\begin{itensq}
\item Determine $q$.
\item Quantos elétrons foram \textbf{retirados} de cada esfera?
\end{itensq}
\dados{$e=\pot{1.6}{-19}\,\si{\coulomb}$}
\resposta{(a) $q=\SI{2}{\micro\coulomb}$; (b) $n=\pot{1.25}{13}$ elétrons}
\resolucao{(a) De $F=\ke q^{2}/r^{2}$,
\[
q=\sqrt{\frac{Fr^{2}}{\ke}}=\sqrt{\frac{0{,}90\cdot(0{,}20)^{2}}{\pot{9}{9}}}
 =\sqrt{\pot{4}{-12}}=\pot{2}{-6}\,\si{\coulomb}.
\]
(b) Repulsão com cargas iguais e sinal positivo indica \emph{falta} de
elétrons: $n=q/e=\pot{2}{-6}/\pot{1.6}{-19}=\pot{1.25}{13}$.}
\end{questao}

\begin{questao}[1]{}
Duas cargas de $\SI{5}{\micro\coulomb}$ cada devem exercer uma sobre a outra
força de $\SI{2.5}{\newton}$ no vácuo. A que distância devem ser colocadas?
\resposta{$r=\SI{0.30}{\meter}$}
\resolucao{Isolando $r$ na lei de Coulomb:
\[
r=\sqrt{\frac{\ke q_1q_2}{F}}
 =\sqrt{\frac{\pot{9}{9}\cdot(\pot{5}{-6})^{2}}{2{,}5}}
 =\sqrt{\frac{0{,}225}{2{,}5}}=\sqrt{0{,}09}=\SI{0.30}{\meter}.
\]}
\end{questao}

\begin{questao}[1]{}
Duas cargas puntiformes se atraem com força $F$. Mantendo as cargas e
\textbf{dobrando} a distância entre elas, a nova força será:
\begin{alternativas}[3]
\item $F/4$
\item $F/2$
\item $F$
\item $2F$
\item $4F$
\end{alternativas}
\resposta{(a)}
\resolucao{$F\propto 1/r^{2}$. Dobrar $r$ divide a força por $2^{2}=4$:
$F'=F/4$. Note que a natureza (atrativa) não muda --- ela depende apenas dos
sinais, nunca da distância.}
\end{questao}

\begin{questao}[1]{}
Duas cargas separadas por uma distância fixa exercem força de
$\SI{0.80}{\newton}$ no vácuo. Todo o conjunto é então imerso em um líquido
isolante de permissividade relativa $\varepsilon_r=4$. Determine a nova força.
\resposta{$F'=\SI{0.20}{\newton}$}
\resolucao{Em um meio dielétrico, $F_{\text{meio}}=F_{\text{vácuo}}/\varepsilon_r$,
pois $\ke$ é substituída por $\ke/\varepsilon_r$:
\[
F'=\frac{0{,}80}{4}=\SI{0.20}{\newton}.
\]
O dielétrico se polariza e blinda parcialmente as cargas --- por isso a força
\emph{sempre} diminui ($\varepsilon_r\ge 1$).}
\end{questao}

\begin{questao}[1]{}
Um bastão carregado positivamente é aproximado --- sem tocar --- de uma esfera
metálica \textbf{neutra} e isolada. A força entre eles é:
\begin{alternativas}[2]
\item nula, pois a esfera é neutra
\item atrativa
\item repulsiva
\item atrativa ou repulsiva, dependendo da distância
\end{alternativas}
\resposta{(b)}
\resolucao{A carga \emph{líquida} da esfera continua nula, mas o bastão
positivo atrai elétrons livres para a face próxima e repele carga positiva
para a face distante (\textbf{indução eletrostática}). A face próxima, negativa,
está \emph{mais perto} do bastão que a face distante, positiva. Como
$F\propto 1/r^{2}$, a atração vence: resultante \textbf{atrativa}. É o mesmo
princípio pelo qual um canudo atritado atrai pedaços de papel neutros.}
\end{questao}

\begin{questao}[1]{}
Estime a força entre duas cargas de $\SI{1}{\coulomb}$ separadas por
$\SI{1}{\kilo\meter}$ e comente o resultado.
\resposta{$F=\SI{9}{\kilo\newton}$ (o peso de cerca de \SI{900}{\kilogram})}
\resolucao{
\[
F=\frac{\pot{9}{9}\cdot 1\cdot 1}{(10^{3})^{2}}=\pot{9}{3}\,\si{\newton}.
\]
Nove mil newtons --- o peso de quase uma tonelada --- a \emph{um quilômetro}
de distância. Isso mostra por que o coulomb é uma unidade enorme para a
eletrostática: cargas realmente separáveis por atrito ficam na faixa de
$\si{\nano\coulomb}$ a $\si{\micro\coulomb}$.}
\end{questao}

\blocoaplicacao

\begin{questao}[2]{}
Três cargas estão fixas sobre o eixo $x$: $q_1=\SI{+3}{\micro\coulomb}$ em
$x=0$, $q_2=\SI{+2}{\micro\coulomb}$ em $x=\SI{20}{\centi\meter}$ e
$q_3=\SI{-5}{\micro\coulomb}$ em $x=\SI{50}{\centi\meter}$. Determine a força
resultante sobre $q_2$ (módulo e sentido).
\resposta{$F=\SI{2.35}{\newton}$, no sentido $+x$}
\resolucao{Trate cada par separadamente e some \emph{com sinal}.
\[
F_{12}=\frac{\pot{9}{9}\cdot\pot{3}{-6}\cdot\pot{2}{-6}}{(0{,}20)^{2}}
=\SI{1.35}{\newton}\quad(\text{repulsão}\Rightarrow +x)
\]
\[
F_{32}=\frac{\pot{9}{9}\cdot\pot{2}{-6}\cdot\pot{5}{-6}}{(0{,}30)^{2}}
=\SI{1.00}{\newton}\quad(\text{atração por } q_3 \Rightarrow +x)
\]
As duas apontam no mesmo sentido: $F=1{,}35+1{,}00=\SI{2.35}{\newton}$ em $+x$.
\textbf{Cuidado:} não some as cargas antes; a lei de Coulomb vale par a par
(princípio da superposição).}
\end{questao}

\begin{questao}[2]{}
Duas esferas metálicas idênticas com $q_1=\SI{+8}{\micro\coulomb}$ e
$q_2=\SI{-2}{\micro\coulomb}$ estão a $\SI{30}{\centi\meter}$. Elas são postas
em contato e recolocadas na mesma posição.
\begin{itensq}
\item Calcule a força antes do contato.
\item Calcule a força depois.
\end{itensq}
\resposta{(a) $\SI{1.6}{\newton}$, atrativa; (b) $\SI{0.9}{\newton}$, repulsiva}
\resolucao{(a) $F=\dfrac{\pot{9}{9}\cdot\pot{8}{-6}\cdot\pot{2}{-6}}{(0{,}30)^{2}}
=\dfrac{0{,}144}{0{,}09}=\SI{1.6}{\newton}$, atrativa (sinais opostos).\\
(b) Esferas \emph{idênticas} em contato dividem a carga total igualmente:
\[
q'=\frac{q_1+q_2}{2}=\frac{+8-2}{2}=\SI{+3}{\micro\coulomb}\ \text{cada}.
\]
\[
F'=\frac{\pot{9}{9}\cdot(\pot{3}{-6})^{2}}{0{,}09}=\SI{0.9}{\newton},
\ \text{agora repulsiva}.
\]
A força \emph{e} a natureza mudaram: o contato conserva a carga total, não a
individual.}
\end{questao}

\begin{questao}[2]{}
Duas cargas fixas exercem entre si $\SI{0.60}{\newton}$ no vácuo. O conjunto é
mergulhado em óleo com $\varepsilon_r=2{,}5$.
\begin{itensq}
\item Qual é a nova força, mantida a distância?
\item Que fração da distância original restauraria a força de
      $\SI{0.60}{\newton}$?
\end{itensq}
\resposta{(a) $\SI{0.24}{\newton}$; (b) $r'=r/\sqrt{2{,}5}\approx 0{,}63\,r$}
\resolucao{(a) $F'=F/\varepsilon_r=0{,}60/2{,}5=\SI{0.24}{\newton}$.\\
(b) Queremos $\dfrac{\ke q_1q_2}{\varepsilon_r r'^{2}}=\dfrac{\ke q_1q_2}{r^{2}}$,
ou seja $\varepsilon_r r'^{2}=r^{2}$:
\[
r'=\frac{r}{\sqrt{\varepsilon_r}}=\frac{r}{\sqrt{2{,}5}}=0{,}632\,r.
\]
Aproximar as cargas em cerca de $37\%$ compensa exatamente o efeito do meio.}
\end{questao}

\begin{questao}[2]{}
A força entre duas cargas $q_1$ e $q_2$ separadas por $d$ vale $F$. Se $q_1$ for
\textbf{triplicada} e a distância for reduzida à \textbf{metade}, a nova força
será:
\begin{alternativas}[3]
\item $1{,}5F$
\item $3F$
\item $6F$
\item $12F$
\item $24F$
\end{alternativas}
\resposta{(d)}
\resolucao{Os efeitos são multiplicativos e independentes:
\[
F'=\ke\frac{(3q_1)q_2}{(d/2)^{2}}=3\cdot 4\cdot \ke\frac{q_1q_2}{d^{2}}=12F.
\]
Fator $3$ da carga $\times$ fator $4$ da distância $=12$.}
\end{questao}

\begin{questao}[2]{}
Em uma impressora a jato eletrostático, cada gota de tinta tem massa
$\pot{1.0}{-11}\,\si{\kilogram}$ e é carregada pela remoção de
$\pot{2.0}{5}$ elétrons. Duas gotas vizinhas ficam a $\SI{1.0}{\milli\meter}$.
Compare a força elétrica entre elas com o peso de uma gota.
\dados{$e=\pot{1.6}{-19}\,\si{\coulomb}$; $g=\SI{9.8}{\meter\per\second\squared}$}
\resposta{$F_e=\pot{9.2}{-12}\,\si{\newton}$; $P=\pot{9.8}{-11}\,\si{\newton}$;
$F_e/P\approx 0{,}094$}
\resolucao{Carga da gota: $q=ne=\pot{2.0}{5}\cdot\pot{1.6}{-19}
=\pot{3.2}{-14}\,\si{\coulomb}$.
\[
F_e=\frac{\pot{9}{9}\cdot(\pot{3.2}{-14})^{2}}{(\pot{1.0}{-3})^{2}}
=\pot{9.2}{-12}\,\si{\newton}.
\]
Peso: $P=mg=\pot{1.0}{-11}\cdot 9{,}8=\pot{9.8}{-11}\,\si{\newton}$.
Razão $\approx 9{,}4\%$: a repulsão entre gotas é pequena diante do peso, mas
não desprezível --- é ela que degrada a nitidez quando as gotas são emitidas
muito próximas.}
\end{questao}

\begin{questao}[2]{}
Duas cargas positivas, $q_1=\SI{4}{\micro\coulomb}$ em $x=0$ e
$q_2=\SI{9}{\micro\coulomb}$ em $x=\SI{50}{\centi\meter}$, estão fixas.
Determine o ponto onde uma terceira carga ficaria em equilíbrio.
\resposta{$x=\SI{20}{\centi\meter}$}
\resolucao{Como as duas são positivas, o ponto de força nula está
\emph{entre} elas. Seja $x$ a distância a $q_1$:
\[
\frac{\ke q_1 q_3}{x^{2}}=\frac{\ke q_2 q_3}{(0{,}50-x)^{2}}
\;\Longrightarrow\;
\frac{4}{x^{2}}=\frac{9}{(0{,}50-x)^{2}}.
\]
Extraindo a raiz (ambos os lados positivos):
$\dfrac{2}{x}=\dfrac{3}{0{,}50-x}\Rightarrow 2(0{,}50-x)=3x
\Rightarrow x=\SI{0.20}{\meter}$.
O ponto fica \emph{mais próximo da carga menor}, como esperado. Note que $q_3$
se cancela: a posição não depende dela.}
\end{questao}

\begin{questao}[2]{}
Uma carga $q$ e outra $100q$ interagem. Sobre os módulos das forças que uma
exerce sobre a outra, é correto afirmar:
\begin{alternativas}[2]
\item a força sobre $q$ é 100 vezes maior
\item a força sobre $100q$ é 100 vezes maior
\item as forças têm o mesmo módulo
\item depende das massas
\end{alternativas}
\resposta{(c)}
\resolucao{A lei de Coulomb é simétrica no produto $q_1q_2$: ambas valem
$\ke\,q\cdot 100q/r^{2}$. Formam um par ação--reação (3\textsuperscript{a} lei
de Newton): mesmo módulo, mesma direção, sentidos opostos. O que \emph{difere}
é a \textbf{aceleração}, se as massas forem diferentes --- confundir força com
aceleração é o erro clássico aqui.}
\end{questao}

\begin{questao}[2]{}
No modelo de Bohr para o hidrogênio, o elétron orbita o próton a
$\pot{5.3}{-11}\,\si{\meter}$. Determine a força elétrica e a aceleração
centrípeta do elétron.
\dados{$e=\pot{1.6}{-19}\,\si{\coulomb}$; $m_e=\pot{9.11}{-31}\,\si{\kilogram}$}
\resposta{$F=\pot{8.2}{-8}\,\si{\newton}$;
$a=\pot{9.0}{22}\,\si{\meter\per\second\squared}$}
\resolucao{
\[
F=\frac{\pot{9}{9}\cdot(\pot{1.6}{-19})^{2}}{(\pot{5.3}{-11})^{2}}
=\frac{\pot{2.304}{-28}}{\pot{2.81}{-21}}=\pot{8.2}{-8}\,\si{\newton}.
\]
\[
a=\frac{F}{m_e}=\frac{\pot{8.2}{-8}}{\pot{9.11}{-31}}
=\pot{9.0}{22}\,\si{\meter\per\second\squared}.
\]
Uma força minúscula produz aceleração de $10^{22}\,\si{\meter\per\second\squared}$
porque a massa do elétron é ainda mais minúscula.}
\end{questao}

\begin{questao}[2]{}
Duas esferas metálicas idênticas com cargas $+Q$ e $+3Q$, separadas por $d$,
repelem-se com força $F$. Elas são postas em contato e depois afastadas até
$2d$. Determine a nova força em função de $F$.
\resposta{$F'=F/3$}
\resolucao{Antes: $F=\ke\dfrac{Q\cdot 3Q}{d^{2}}=\dfrac{3\ke Q^{2}}{d^{2}}$.\\
Após o contato, cada esfera fica com $\dfrac{Q+3Q}{2}=2Q$. A $2d$:
\[
F'=\ke\frac{(2Q)(2Q)}{(2d)^{2}}=\ke\frac{4Q^{2}}{4d^{2}}=\frac{\ke Q^{2}}{d^{2}}
=\frac{F}{3}.
\]
O produto das cargas subiu (de $3Q^2$ para $4Q^2$), mas o fator $4$ da
distância dominou.}
\end{questao}

\begin{questao}[2]{}
Compare a força entre duas cargas de $\SI{1.0}{\micro\coulomb}$ separadas por
$\SI{1.0}{\centi\meter}$ com o peso de um corpo de $\SI{1.0}{\kilogram}$, e
interprete.
\dados{$g=\SI{9.8}{\meter\per\second\squared}$}
\resposta{$F=\SI{90}{\newton}$ contra $P=\SI{9.8}{\newton}$; razão $\approx 9{,}2$}
\resolucao{
\[
F=\frac{\pot{9}{9}\cdot(\pot{1}{-6})^{2}}{(\pot{1}{-2})^{2}}
=\frac{\pot{9}{-3}}{\pot{1}{-4}}=\SI{90}{\newton};
\qquad P=1{,}0\cdot 9{,}8=\SI{9.8}{\newton}.
\]
Um milionésimo de coulomb, a um centímetro, produz nove vezes o peso de um
quilograma. Essa desproporção entre a força elétrica e a gravitacional é a razão
pela qual é a matéria macroscópica é praticamente neutra: qualquer desequilíbrio
apreciável de carga seria violentamente desfeito.}
\end{questao}

\blocoaprofundamento

\begin{questao}[3]{}
Duas cargas fixas sobre o eixo $x$: $q_1=\SI{+4}{\micro\coulomb}$ em $x=0$ e
$q_2=\SI{-1}{\micro\coulomb}$ em $x=\SI{30}{\centi\meter}$. Determine todos os
pontos do eixo onde a força sobre uma carga de prova seria nula.
\resposta{Um único ponto: $x=\SI{60}{\centi\meter}$}
\resolucao{Com \textbf{sinais opostos}, o ponto não pode estar entre as cargas
(ali as duas forças apontam para o mesmo lado). Deve estar \emph{fora} do
segmento e, mais especificamente, do lado da carga de \emph{menor} módulo ---
só assim a maior proximidade compensa o menor valor. Testando $x>0{,}30$:
\[
\frac{\ke\cdot 4\mu}{x^{2}}=\frac{\ke\cdot 1\mu}{(x-0{,}30)^{2}}
\;\Longrightarrow\;
\frac{2}{x}=\frac{1}{x-0{,}30}
\;\Longrightarrow\;
2x-0{,}60=x \Rightarrow x=\SI{0.60}{\meter}.
\]
\textbf{Verificação:}
$\ke\cdot4\mu/(0{,}60)^{2}=\ke\cdot1\mu/(0{,}30)^{2}$ --- confere. A raiz
$x=0{,}20$ da equação quadrática completa é \emph{espúria}: corresponde ao
ponto interno, onde as forças se somam em vez de se cancelar.}
\end{questao}

\begin{questao}[3]{}
Seis cargas iguais $Q=\SI{2}{\micro\coulomb}$ ocupam os vértices de um hexágono
regular de lado $a=\SI{10}{\centi\meter}$, com uma carga
$q=\SI{1}{\micro\coulomb}$ no centro.
\begin{itensq}
\item Qual é a força resultante sobre $q$?
\item Uma das seis cargas é removida. Qual passa a ser a força sobre $q$
      (módulo e direção)?
\end{itensq}
\resposta{(a) nula; (b) $\SI{1.8}{\newton}$, apontando \emph{para} o vértice vazio}
\resolucao{(a) Em um hexágono regular, o centro dista $a$ de todos os vértices
e as cargas ocupam posições diametralmente opostas duas a duas. Cada par
contribui com forças de mesmo módulo e sentidos opostos: resultante nula, por
simetria.\\
(b) Use o argumento da superposição \emph{ao contrário}: 
$\vec{F}_{6}= \vec{F}_{5}+\vec{F}_{\text{removida}}=\vec{0}$, logo
$\vec{F}_{5}=-\vec{F}_{\text{removida}}$.
\[
|\vec{F}_{5}|=\frac{\ke Qq}{a^{2}}
=\frac{\pot{9}{9}\cdot\pot{2}{-6}\cdot\pot{1}{-6}}{(0{,}10)^{2}}
=\SI{1.8}{\newton}.
\]
A carga removida \emph{repelia} $q$ para longe do vértice; sem ela, a resultante
aponta \textbf{para} o vértice vazio. Este truque de simetria evita somar cinco
vetores.}
\end{questao}

\begin{questao}[3]{}
Três cargas ocupam os vértices de um triângulo retângulo:
$q_A=\SI{+2}{\micro\coulomb}$ na origem, $q_B=\SI{+3}{\micro\coulomb}$ em
$(\SI{30}{\centi\meter},0)$ e $q_C=\SI{+4}{\micro\coulomb}$ em
$(0,\SI{40}{\centi\meter})$. Determine a força resultante sobre $q_A$: módulo e
ângulo com o eixo $x$.
\resposta{$F=\SI{0.75}{\newton}$, a $\SI{216.9}{\degree}$ (isto é, $\SI{36.9}{\degree}$ abaixo
do semieixo $-x$)}
\resolucao{As duas forças sobre $q_A$ são perpendiculares --- o que torna a
composição imediata.
\[
F_{AB}=\frac{\pot{9}{9}\cdot\pot{2}{-6}\cdot\pot{3}{-6}}{(0{,}30)^{2}}
=\SI{0.60}{\newton}\ \ (\text{repulsão}\Rightarrow -x)
\]
\[
F_{AC}=\frac{\pot{9}{9}\cdot\pot{2}{-6}\cdot\pot{4}{-6}}{(0{,}40)^{2}}
=\SI{0.45}{\newton}\ \ (\text{repulsão}\Rightarrow -y)
\]
\[
F=\sqrt{0{,}60^{2}+0{,}45^{2}}=\sqrt{0{,}5625}=\SI{0.75}{\newton},
\qquad
\tan\theta=\frac{0{,}45}{0{,}60}=0{,}75 \Rightarrow \theta=\SI{36.87}{\degree}.
\]
Um clássico triângulo $3$--$4$--$5$ escondido. \textbf{Atenção:} a hipotenusa
$BC=\SI{50}{\centi\meter}$ é irrelevante --- ela governa a força \emph{entre}
$q_B$ e $q_C$, não a força sobre $q_A$.}
\end{questao}

\begin{questao}[3]{}
Duas pequenas esferas metálicas idênticas, separadas por $\SI{10}{\centi\meter}$,
\textbf{atraem-se} com força de $\SI{0.108}{\newton}$. Postas em contato e
recolocadas na mesma distância, passam a \textbf{repelir-se} com
$\SI{0.036}{\newton}$. Determine as cargas iniciais.
\resposta{$q_1=\SI{+0.6}{\micro\coulomb}$ e $q_2=\SI{-0.2}{\micro\coulomb}$
(ou a permuta)}
\resolucao{\textbf{Etapa 1 --- produto.} A atração indica sinais opostos:
\[
|q_1q_2|=\frac{F_1r^{2}}{\ke}=\frac{0{,}108\cdot 0{,}01}{\pot{9}{9}}
=\pot{1.2}{-13}\ \Rightarrow\ q_1q_2=-\pot{1.2}{-13}\,\si{\coulomb\squared}.
\]
\textbf{Etapa 2 --- soma.} Após o contato cada esfera fica com
$q'=(q_1+q_2)/2$:
\[
q'=\sqrt{\frac{F_2r^{2}}{\ke}}=\sqrt{\frac{0{,}036\cdot0{,}01}{\pot{9}{9}}}
=\pot{2}{-7}\,\si{\coulomb}
\ \Rightarrow\ q_1+q_2=\pot{4}{-7}\,\si{\coulomb}.
\]
\textbf{Etapa 3 --- soma e produto.} As cargas são raízes de
$x^{2}-Sx+P=0$ com $S=\pot{4}{-7}$ e $P=-\pot{1.2}{-13}$:
\[
x=\frac{\pot{4}{-7}\pm\sqrt{\pot{1.6}{-13}+\pot{4.8}{-13}}}{2}
 =\frac{\pot{4}{-7}\pm\pot{8}{-7}}{2}
 \Rightarrow x=\pot{6}{-7}\ \text{ou}\ -\pot{2}{-7}.
\]
A repulsão final confirma que a soma é não nula. Se as cargas fossem simétricas
($q$ e $-q$), o contato as neutralizaria e a força final seria zero.}
\end{questao}

\begin{questao}[3]{}
Um pequeno bloco de massa $m=\SI{20}{\gram}$, com carga $q=\SI{1}{\micro\coulomb}$,
repousa sobre um plano inclinado \textbf{liso} de $\SI{30}{\degree}$. Uma carga fixa
$Q=\SI{1}{\micro\coulomb}$ está presa na base do plano, e a força elétrica atua
ao longo da superfície. Determine a distância de equilíbrio $d$ medida sobre o
plano.
\dados{$g=\SI{9.8}{\meter\per\second\squared}$}
\resposta{$d\approx\SI{30.3}{\centi\meter}$}
\resolucao{Sobre o plano liso, o equilíbrio na direção da superfície envolve
apenas a componente do peso e a repulsão elétrica (a normal é perpendicular):
\[
\frac{\ke Qq}{d^{2}}=mg\sin\theta.
\]
\[
mg\sin\SI{30}{\degree}=0{,}020\cdot 9{,}8\cdot 0{,}5=\SI{0.098}{\newton};
\qquad
d=\sqrt{\frac{\pot{9}{9}\cdot(\pot{1}{-6})^{2}}{0{,}098}}
=\sqrt{0{,}0918}=\SI{0.303}{\meter}.
\]
\textbf{Estabilidade:} se o bloco descer, $d$ diminui e a repulsão cresce como
$1/d^{2}$, empurrando-o de volta; se subir, a repulsão cai e o peso o traz de
volta. O equilíbrio é \emph{estável}.}
\end{questao}

\begin{questao}[3]{}
Um dipolo é formado por $+q$ e $-q$, com $q=\SI{1}{\micro\coulomb}$, fixos em
$(\pm\SI{5}{\centi\meter},0)$. Uma carga $Q=\SI{2}{\micro\coulomb}$ é colocada
em $(0,\SI{12}{\centi\meter})$, sobre a mediatriz.
\begin{itensq}
\item Mostre que a resultante sobre $Q$ é paralela ao eixo do dipolo.
\item Calcule seu módulo.
\end{itensq}
\resposta{(a) as componentes verticais se cancelam; (b) $F=\SI{0.82}{\newton}$,
no sentido $-x$ (de $+q$ para $-q$)}
\resolucao{(a) Cada carga dista
$r=\sqrt{0{,}05^{2}+0{,}12^{2}}=\SI{0.13}{\meter}$ de $Q$ (triângulo
$5$--$12$--$13$). A carga $+q$ \emph{repele} $Q$ e $-q$ a \emph{atrai}: as duas
forças têm o mesmo módulo e as componentes em $y$ apontam em sentidos opostos,
cancelando-se. Restam as componentes em $x$, que se \textbf{somam}.\\
(b) Cada força vale $F_1=\ke qQ/r^{2}$ e sua componente horizontal é
$F_1\cdot(a/r)$, com $a=\SI{0.05}{\meter}$:
\[
F=2\,\frac{\ke qQ}{r^{2}}\cdot\frac{a}{r}=\frac{2\ke qQ\,a}{r^{3}}
=\frac{2\cdot\pot{9}{9}\cdot\pot{1}{-6}\cdot\pot{2}{-6}\cdot 0{,}05}
       {(0{,}13)^{3}}=\SI{0.82}{\newton}.
\]
Note a dependência $1/r^{3}$: longe do dipolo a força cai \emph{mais rápido}
que a de uma carga isolada, porque as duas cargas quase se cancelam.}
\end{questao}

\blocodesafio

\begin{questao}[4]{}
\textbf{(Modelagem experimental.)} Em uma reprodução da balança de torção, duas
esferas com cargas iguais $q=\SI{0.10}{\micro\coulomb}$ foram mantidas a várias
distâncias e a força foi medida:

\begin{center}
\begin{tabular}{cccccc}
\hline
$r$ (m) & 0,020 & 0,030 & 0,040 & 0,050 & 0,060\\
$F$ (N) & 0,223 & 0,101 & 0,0558 & 0,0362 & 0,0249\\
\hline
\end{tabular}
\end{center}

\begin{itensq}
\item Proponha uma linearização que permita testar a hipótese $F=C\,r^{n}$.
\item Estime $n$ e discuta se os dados sustentam a lei do inverso do quadrado.
\item Obtenha o valor experimental de $\ke$ e compare com o tabelado.
\end{itensq}
\resposta{(a) $\log F=\log C+n\log r$; (b) $n\approx-2{,}00$; (c)
$\ke\approx\pot{9.1}{9}\,\si{\newton\meter\squared\per\coulomb\squared}$
(erro $\approx 1\%$)}
\resolucao{(a) Tomando logaritmos, uma lei de potência vira uma reta:
$\log F = \log C + n\log r$. O coeficiente angular do gráfico
$\log F \times \log r$ é o expoente procurado.\\
(b) Ajustando os cinco pontos por mínimos quadrados obtém-se
$n=-1{,}997$ e $C=\pot{9.09}{-5}$. O expoente reproduz $-2$ com desvio de
$0{,}15\%$ --- dentro do que se espera de medidas de força dessa ordem. Um teste
rápido sem regressão: dobrar $r$ de $0{,}020$ para $0{,}040$ deveria dividir $F$
por $4$; de fato $0{,}223/0{,}0558=4{,}00$.\\
(c) Como $C=\ke q^{2}$ e $q^{2}=\pot{1}{-14}\,\si{\coulomb\squared}$:
\[
\ke=\frac{C}{q^{2}}=\frac{\pot{9.09}{-5}}{\pot{1}{-14}}
=\pot{9.09}{9}\,\si{\newton\meter\squared\per\coulomb\squared},
\]
contra $\pot{8.99}{9}$ tabelado --- erro relativo de $1{,}1\%$.\\
\textbf{Observação de método:} a linearização logarítmica é preferível a ajustar
$F\times 1/r^{2}$ diretamente, porque \emph{não pressupõe} o expoente --- ela o
mede. Esse é exatamente o argumento que Coulomb precisou construir em 1785.}
\end{questao}

\begin{questao}[4]{}
\textbf{(Otimização com cálculo.)} Um anel fino de raio $R$, com carga total $Q$
uniformemente distribuída, está no plano $yz$. Uma carga puntiforme $q$ é
colocada sobre o eixo do anel, a uma distância $x$ do centro.
\begin{itensq}
\item Mostre que $F(x)=\dfrac{\ke Q q\,x}{(x^{2}+R^{2})^{3/2}}$.
\item Determine o valor de $x$ que \textbf{maximiza} a força.
\item Calcule $F_{\max}$ para $Q=\SI{10}{\micro\coulomb}$,
      $q=\SI{2}{\micro\coulomb}$ e $R=\SI{10}{\centi\meter}$.
\end{itensq}
\resposta{(b) $x=R/\sqrt{2}\approx\SI{7.07}{\centi\meter}$;
(c) $F_{\max}=\dfrac{2\ke Qq}{3\sqrt{3}\,R^{2}}=\SI{6.93}{\newton}$}
\resolucao{(a) Divida o anel em elementos $\mathrm{d} Q$. Cada um dista
$s=\sqrt{x^{2}+R^{2}}$ de $q$ e contribui com $\ke q\,\mathrm{d} Q/s^{2}$. Por
simetria, as componentes perpendiculares ao eixo se cancelam aos pares; sobra a
projeção axial, fator $\cos\alpha=x/s$:
\[
F=\int \frac{\ke q\,\mathrm{d} Q}{s^{2}}\cdot\frac{x}{s}
 =\frac{\ke q x}{(x^{2}+R^{2})^{3/2}}\int \mathrm{d} Q
 =\frac{\ke Q q\,x}{(x^{2}+R^{2})^{3/2}}.
\]
(b) Derivando e igualando a zero:
\[
\frac{\mathrm{d} F}{\mathrm{d} x}\propto (x^{2}+R^{2})^{3/2}-x\cdot 3x(x^{2}+R^{2})^{1/2}=0
\;\Rightarrow\; x^{2}+R^{2}-3x^{2}=0 \;\Rightarrow\; x=\frac{R}{\sqrt{2}}.
\]
(c) Substituindo $x=R/\sqrt2$:
\[
F_{\max}=\frac{2\,\ke Qq}{3\sqrt{3}\,R^{2}}
=\frac{2\cdot\pot{9}{9}\cdot\pot{1}{-5}\cdot\pot{2}{-6}}{3\sqrt3\cdot 0{,}01}
=\frac{0{,}36}{0{,}0520}=\SI{6.93}{\newton}.
\]
\textbf{Casos-limite:} em $x=0$ a força é nula (simetria perfeita); para
$x\gg R$, $F\to\ke Qq/x^{2}$ --- o anel se comporta como carga puntiforme.
Entre esses dois zeros deve existir um máximo, e ele está em $R/\sqrt2$.}
\end{questao}

\begin{questao}[4]{}
\textbf{(Distribuição contínua.)} Um bastão isolante retilíneo de comprimento
$L$ possui carga $Q$ uniformemente distribuída. Uma carga puntiforme $q$ está
sobre o prolongamento do bastão, a uma distância $d$ da extremidade mais
próxima.
\begin{itensq}
\item Deduza $F=\dfrac{\ke Qq}{d\,(d+L)}$.
\item Verifique o comportamento para $L\ll d$.
\item Calcule $F$ para $Q=\SI{4}{\micro\coulomb}$, $L=\SI{20}{\centi\meter}$,
      $q=\SI{1}{\micro\coulomb}$ e $d=\SI{10}{\centi\meter}$.
\end{itensq}
\resposta{(b) $F\to\ke Qq/d^{2}$; (c) $F=\SI{1.2}{\newton}$}
\resolucao{(a) Densidade linear $\lambda=Q/L$. Um elemento $\mathrm{d} x$ a uma
distância $x$ da carga carrega $\mathrm{d} Q=\lambda\,\mathrm{d} x$ e todas as contribuições
são \emph{colineares} --- por isso a integral é escalar:
\[
F=\int_{d}^{d+L}\frac{\ke q\lambda\,\mathrm{d} x}{x^{2}}
 =\ke q\lambda\left[-\frac{1}{x}\right]_{d}^{d+L}
 =\ke q\lambda\left(\frac{1}{d}-\frac{1}{d+L}\right)
 =\frac{\ke q\lambda L}{d(d+L)}=\frac{\ke Qq}{d(d+L)}.
\]
(b) Se $L\ll d$, então $d+L\approx d$ e $F\to\ke Qq/d^{2}$: recuperamos a lei de
Coulomb puntiforme, como deve ser.\\
(c) $F=\dfrac{\pot{9}{9}\cdot\pot{4}{-6}\cdot\pot{1}{-6}}{0{,}10\cdot 0{,}30}
=\dfrac{0{,}036}{0{,}03}=\SI{1.2}{\newton}$.\\
\textbf{Compare:} se toda a carga estivesse concentrada no centro do bastão
($\SI{20}{\centi\meter}$), teríamos $\SI{0.9}{\newton}$. A força real é maior
porque as porções mais próximas pesam com $1/x^{2}$ --- concentrar carga no
centro geométrico \emph{subestima} a força.}
\end{questao}

\begin{questao}[4]{}
\textbf{(Demonstração --- estabilidade.)} Duas cargas $+Q$ estão fixas em
$x=-d$ e $x=+d$. Uma carga $q$ é colocada na origem.
\begin{itensq}
\item Mostre que a origem é posição de equilíbrio para qualquer $q$.
\item Para $q>0$, analise a estabilidade a um pequeno deslocamento
      \emph{longitudinal} ($\pm x$).
\item Repita para um deslocamento \emph{transversal} ($\pm y$) e conclua.
\end{itensq}
\resposta{(a) por simetria; (b) estável, com $F\approx-4\ke Qq\,x/d^{3}$;
(c) instável, com $F\approx+2\ke Qq\,y/d^{3}$ --- não há equilíbrio estável em
todas as direções}
\resolucao{(a) As duas forças têm o mesmo módulo e sentidos opostos: resultante
nula, independentemente do sinal e do valor de $q$.\\
(b) Deslocando $q>0$ para $x$ pequeno:
\[
F_x=\ke Qq\left[\frac{1}{(d+x)^{2}}-\frac{1}{(d-x)^{2}}\right].
\]
Com $(d\pm x)^{-2}\approx d^{-2}(1\mp 2x/d)$:
\[
F_x\approx\frac{\ke Qq}{d^{2}}\left(-\frac{4x}{d}\right)=-\frac{4\ke Qq}{d^{3}}\,x.
\]
Sinal negativo $\Rightarrow$ força \textbf{restauradora}: estável nessa direção
(e o movimento é harmônico simples).\\
(c) Deslocando para $y$ pequeno, cada carga dista $\sqrt{d^{2}+y^{2}}\approx d$ e
repele $q$ com componente transversal $\ke Qq\,y/d^{3}$; as duas se somam:
\[
F_y\approx +\frac{2\ke Qq}{d^{3}}\,y,
\]
sinal positivo $\Rightarrow$ força que \textbf{afasta}: instável.\\
\textbf{Conclusão.} O equilíbrio é do tipo \emph{ponto de sela}: estável em uma
direção, instável na outra. Este é um caso particular do \textbf{teorema de
Earnshaw}: nenhuma configuração de cargas fixas pode manter uma carga puntiforme
em equilíbrio estável apenas por forças eletrostáticas. É por isso que modelos
puramente eletrostáticos do átomo não se sustentam.}
\end{questao}

\begin{questao}[4]{}
\textbf{(Série infinita.)} Cargas puntiformes de mesmo módulo $Q$ e
\textbf{sinais alternados} ocupam as posições $x=a,\,2a,\,3a,\dots$ sobre um
eixo: $-Q$ em $a$, $+Q$ em $2a$, $-Q$ em $3a$ e assim por diante,
indefinidamente. Uma carga $+q$ está na origem.
\begin{itensq}
\item Escreva a força resultante sobre $q$ como uma série.
\item Some a série e obtenha a expressão fechada.
\end{itensq}
\dados{$\displaystyle\sum_{n=1}^{\infty}\frac{(-1)^{n+1}}{n^{2}}=\frac{\pi^{2}}{12}$}
\resposta{$F=\dfrac{\pi^{2}}{12}\cdot\dfrac{\ke Qq}{a^{2}}
\approx 0{,}822\,\dfrac{\ke Qq}{a^{2}}$, atrativa (sentido $+x$)}
\resolucao{(a) A $n$-ésima carga está em $x=na$ e tem sinal $(-1)^{n}Q$. Uma
carga negativa \emph{atrai} $q$ (puxa no sentido $+x$); uma positiva
\emph{repele} (sentido $-x$). Tomando $+x$ como positivo:
\[
F=\frac{\ke Qq}{a^{2}}\left(\frac{1}{1^{2}}-\frac{1}{2^{2}}+\frac{1}{3^{2}}
-\frac{1}{4^{2}}+\cdots\right)
=\frac{\ke Qq}{a^{2}}\sum_{n=1}^{\infty}\frac{(-1)^{n+1}}{n^{2}}.
\]
(b) Usando o resultado dado:
\[
F=\frac{\pi^{2}}{12}\cdot\frac{\ke Qq}{a^{2}}\approx 0{,}822\,\frac{\ke Qq}{a^{2}}.
\]
\textbf{Interpretação.} Apesar de haver infinitas cargas, a força é
\emph{finita} e vale menos que a da primeira carga sozinha: a alternância de
sinais produz cancelamento progressivo, e o decaimento $1/n^{2}$ garante
convergência absoluta. Se todas as cargas fossem negativas, a soma seria
$\pi^2/6\approx1{,}645$ --- exatamente o dobro.}
\end{questao}

\begin{questao}[4]{}
\textbf{(Oscilações.)} Uma partícula de massa $m=\SI{1.0}{\gram}$ e carga
$-q=\SI{-2}{\micro\coulomb}$ é liberada em repouso sobre o eixo de um anel de
raio $R=\SI{10}{\centi\meter}$ e carga $Q=\SI{5}{\micro\coulomb}$, a uma
distância $x_0\ll R$ do centro.
\begin{itensq}
\item Mostre que, para $x\ll R$, o movimento é harmônico simples.
\item Obtenha $\omega$ e o período $T$.
\end{itensq}
\resposta{(b) $\omega=\sqrt{\ke Qq/(mR^{3})}=\SI{300}{\radian\per\second}$;
$T\approx\SI{20.9}{\milli\second}$}
\resolucao{(a) A força axial sobre a carga negativa é atrativa, dirigida ao
centro, de módulo $F=\ke Qq\,x/(x^{2}+R^{2})^{3/2}$. Para $x\ll R$, o
denominador tende a $R^{3}$:
\[
F_x\approx-\frac{\ke Qq}{R^{3}}\,x.
\]
Força proporcional ao deslocamento e de sinal oposto --- exatamente a forma
$F=-kx$ da lei de Hooke, com constante efetiva $k_{\text{ef}}=\ke Qq/R^{3}$.
Logo o MHS.\\
(b)
\[
\omega=\sqrt{\frac{k_{\text{ef}}}{m}}=\sqrt{\frac{\ke Qq}{mR^{3}}}
=\sqrt{\frac{\pot{9}{9}\cdot\pot{5}{-6}\cdot\pot{2}{-6}}
{\pot{1}{-3}\cdot(0{,}10)^{3}}}
=\sqrt{\frac{0{,}09}{\pot{1}{-6}}}=\SI{300}{\radian\per\second},
\]
\[
T=\frac{2\pi}{\omega}=\SI{20.9}{\milli\second}.
\]
\textbf{Observe:} o período não depende de $x_0$ --- é a assinatura do MHS. Fora
da aproximação $x\ll R$, a força deixa de ser linear e o movimento continua
periódico, mas com período dependente da amplitude.}
\end{questao}

\begin{questao}[4]{}
\textbf{(Modelagem com taxas.)} Duas esferas idênticas de massa
$m=\SI{2.0}{\gram}$, cada uma com carga $q$, pendem de fios isolantes de
comprimento $L=\SI{60}{\centi\meter}$ presos ao mesmo ponto. Elas se afastam até
uma separação $x\ll L$. Por fuga de carga para o ar úmido, as esferas se
aproximam lentamente.
\begin{itensq}
\item Mostre que $x^{3}=\dfrac{2\ke q^{2}L}{mg}$.
\item Demonstre que
      $\dfrac{\mathrm{d} x}{\mathrm{d} t}=\dfrac{2x}{3q}\dfrac{\mathrm{d} q}{\mathrm{d} t}$.
\item Se em certo instante $x=\SI{6.0}{\centi\meter}$ e as esferas se aproximam
      a $\SI{1.0}{\milli\meter\per\second}$, calcule $\mathrm{d} q/\mathrm{d} t$.
\end{itensq}
\dados{$g=\SI{9.8}{\meter\per\second\squared}$}
\resposta{(c) $q\approx\SI{19.8}{\nano\coulomb}$ e
$\mathrm{d} q/\mathrm{d} t\approx\SI{-0.49}{\nano\coulomb\per\second}$}
\resolucao{(a) Para cada esfera: tensão, peso e repulsão. Do triângulo de
forças, $\tan\theta=F/(mg)$. Com $x\ll L$,
$\tan\theta\approx\sin\theta=(x/2)/L$. Igualando:
\[
\frac{x}{2L}=\frac{\ke q^{2}}{x^{2}mg}
\;\Longrightarrow\;
x^{3}=\frac{2\ke q^{2}L}{mg}.
\]
(b) Derivando implicitamente em relação ao tempo:
\[
3x^{2}\frac{\mathrm{d} x}{\mathrm{d} t}=\frac{4\ke L}{mg}\,q\,\frac{\mathrm{d} q}{\mathrm{d} t}.
\]
De (a), $\dfrac{2\ke L}{mg}=\dfrac{x^{3}}{q^{2}}$; substituindo:
\[
3x^{2}\frac{\mathrm{d} x}{\mathrm{d} t}=\frac{2x^{3}}{q}\frac{\mathrm{d} q}{\mathrm{d} t}
\;\Longrightarrow\;
\frac{\mathrm{d} x}{\mathrm{d} t}=\frac{2x}{3q}\frac{\mathrm{d} q}{\mathrm{d} t}.
\]
(c) Primeiro $q$, por (a):
\[
q=\sqrt{\frac{x^{3}mg}{2\ke L}}
=\sqrt{\frac{(0{,}060)^{3}\cdot 0{,}0020\cdot 9{,}8}{2\cdot\pot{9}{9}\cdot0{,}60}}
=\pot{1.98}{-8}\,\si{\coulomb}.
\]
Com $\mathrm{d} x/\mathrm{d} t=-\pot{1.0}{-3}\,\si{\meter\per\second}$ (aproximação):
\[
\frac{\mathrm{d} q}{\mathrm{d} t}=\frac{3q}{2x}\frac{\mathrm{d} x}{\mathrm{d} t}
=\frac{3\cdot\pot{1.98}{-8}}{2\cdot 0{,}060}\cdot(-\pot{1}{-3})
=-\pot{4.9}{-10}\,\si{\coulomb\per\second}.
\]
\textbf{Leitura física:} medir uma \emph{velocidade} converte-se em medir uma
\emph{corrente de fuga} de meio nanoampère --- pequena demais para muitos
amperímetros, mas visível a olho nu pelo movimento das esferas.}
\end{questao}

\begin{questao}[4]{}
\textbf{(Simetria e integração.)} Um fio semicircular de raio
$R=\SI{5.0}{\centi\meter}$ possui carga $Q=\SI{3}{\micro\coulomb}$ distribuída
uniformemente. Uma carga $q=\SI{1}{\micro\coulomb}$ ocupa o centro do
semicírculo.
\begin{itensq}
\item Justifique por que a resultante é dirigida ao longo do eixo de simetria.
\item Mostre que $F=\dfrac{2\ke Qq}{\pi R^{2}}$ e calcule seu valor.
\end{itensq}
\resposta{$F=\SI{6.9}{\newton}$, ao longo do eixo de simetria}
\resolucao{(a) Para cada elemento do fio existe outro simétrico em relação ao
eixo vertical. Suas contribuições têm componentes perpendiculares ao eixo iguais
e opostas --- que se cancelam --- e componentes ao longo do eixo que se somam.\\
(b) Densidade linear $\lambda=Q/(\pi R)$. O elemento em $\theta$ tem
$\mathrm{d} Q=\lambda R\,\mathrm{d}\theta$, dista $R$ do centro e contribui com
$\ke q\,\mathrm{d} Q/R^{2}$; sua projeção no eixo vale $\sin\theta$:
\[
F=\int_{0}^{\pi}\frac{\ke q\lambda R\,\mathrm{d}\theta}{R^{2}}\sin\theta
 =\frac{\ke q\lambda}{R}\underbrace{\int_{0}^{\pi}\sin\theta\,\mathrm{d}\theta}_{=2}
 =\frac{2\ke q\lambda}{R}=\frac{2\ke Qq}{\pi R^{2}}.
\]
\[
F=\frac{2\cdot\pot{9}{9}\cdot\pot{3}{-6}\cdot\pot{1}{-6}}
{\pi\,(0{,}050)^{2}}=\frac{0{,}054}{\pot{7.85}{-3}}=\SI{6.9}{\newton}.
\]
\textbf{Compare:} se a mesma carga $Q$ estivesse concentrada em um ponto a $R$,
teríamos $\ke Qq/R^{2}=\SI{10.8}{\newton}$. O fator $2/\pi\approx 0{,}64$ mede
exatamente a perda por cancelamento vetorial ao longo do arco. Um anel
\emph{completo} levaria esse fator a zero.}
\end{questao}

\begin{questao}[4]{}
\textbf{(Síntese --- equilíbrio de um sistema.)} Quatro cargas iguais
$Q=\SI{2}{\micro\coulomb}$ ocupam os vértices de um quadrado de lado $a$. Uma
quinta carga $q$ é colocada no centro.
\begin{itensq}
\item Determine $q$ (sinal e módulo, em função de $Q$) para que \textbf{cada}
      carga dos vértices fique em equilíbrio.
\item Calcule $q$ numericamente.
\item O sistema completo está em equilíbrio estável? Justifique.
\end{itensq}
\resposta{(a) $q=-\dfrac{Q(2\sqrt{2}+1)}{4}$; (b) $q\approx\SI{-1.91}{\micro\coulomb}$;
(c) não --- o equilíbrio é instável}
\resolucao{(a) Por simetria, basta impor o equilíbrio de \emph{um} vértice; a
resultante das outras três cargas aponta ao longo da diagonal, para fora.
\begin{itemize}
\item Duas vizinhas, a distância $a$: cada uma contribui $\ke Q^{2}/a^{2}$, e
suas projeções sobre a diagonal valem $\cos\SI{45}{\degree}=\sqrt2/2$ cada, somando
$\sqrt{2}\,\ke Q^{2}/a^{2}$.
\item A oposta, a distância $a\sqrt{2}$: contribui
$\ke Q^{2}/(2a^{2})$, já ao longo da diagonal.
\end{itemize}
\[
F_{\text{vértices}}=\frac{\ke Q^{2}}{a^{2}}\left(\sqrt{2}+\frac{1}{2}\right).
\]
A carga central dista $a/\sqrt{2}$ do vértice, logo atua com
$F_{c}=\ke|q|Q/(a^{2}/2)=2\ke|q|Q/a^{2}$ e deve ser \textbf{atrativa} ---
portanto $q<0$. Igualando:
\[
2|q|=Q\left(\sqrt2+\frac12\right)
\;\Longrightarrow\;
|q|=\frac{Q\,(2\sqrt{2}+1)}{4}\approx 0{,}957\,Q.
\]
(b) $q\approx-0{,}957\cdot 2=\SI{-1.91}{\micro\coulomb}$.\\
(c) \textbf{Não.} A condição imposta zera a resultante sobre os vértices, mas
não sobre a carga central --- que já está em equilíbrio por simetria, embora
instável. Mais fundamentalmente, o teorema de Earnshaw proíbe equilíbrio
eletrostático estável: qualquer perturbação de $q$ ao longo de uma diagonal a
faz cair sobre um vértice. O equilíbrio existe, mas não sobrevive a uma
perturbação infinitesimal.}
\end{questao}

% END BANCO COMPLEMENTAR

\gabarito[1.2 Lei de Coulomb]

% =====================================================================
\subsection{Campo elétrico}
% =====================================================================

\begin{conceito}[Antes de começar]
\textbf{Campo de uma carga puntiforme:}
\[
E = \ke\,\frac{|q|}{d^2}\quad(\si{\newton\per\coulomb}\ \text{ou}\ \si{\volt\per\meter}).
\]
O campo \emph{aponta para fora} de cargas positivas e \emph{para dentro} de cargas
negativas.\par
\textbf{Força sobre uma carga no campo:} $F = qE$ (mesmo sentido de $E$ se $q>0$;
oposto se $q<0$).\par
\textbf{Superposição:} o campo resultante é a \emph{soma vetorial} dos campos
individuais.\par
\textbf{Ponto de campo nulo} entre duas cargas de mesmo sinal: onde os módulos se
igualam, $\dfrac{|q_1|}{d_1^2} = \dfrac{|q_2|}{d_2^2}$.
\end{conceito}

\legendaniveis

\blocofixacao

\begin{questao}[1]{}
Qual a intensidade do campo elétrico gerado por uma carga de
$\SI{-15}{\micro\coulomb}$, a $\SI{10}{\milli\meter}$ dela, em um meio de constante
$\ke = \SI{8e9}{\newton\meter\squared\per\coulomb\squared}$?
\resposta{$E = \pot{1.2}{9}\,\si{\newton\per\coulomb}$}
\resolucao{Com $d = \SI{10}{\milli\meter} = \SI{0.01}{\meter}$:
\[
E=\ke\frac{|q|}{d^2}
=8\cdot10^9\cdot\frac{\pot{15}{-6}}{(0{,}01)^2}=\pot{1.2}{9}\,\si{\newton\per\coulomb}.
\]
O campo aponta \emph{para} a carga, por ela ser negativa.}
\end{questao}

\begin{questao}[1]{}
O módulo do campo elétrico de uma carga puntiforme $q$, em um ponto P a uma dada
distância, vale $E$. Afastando-se a carga de modo que a distância a P
\textbf{dobre}, o campo em P passa a ser:
\begin{alternativas}[3]
\item $E$
\item $E/2$
\item $E/4$
\item $2E$
\item $4E$
\end{alternativas}
\resposta{(c)}
\resolucao{Como $E \propto \dfrac{1}{d^2}$, dobrar a distância divide o campo por
$2^2 = 4$: o novo valor é $E/4$.}
\end{questao}

\blocoaplicacao

\begin{questao}[2]{}
Uma carga de prova de $\SI{e-9}{\coulomb}$, colocada em um ponto P de um campo
elétrico, fica sujeita a uma força de $\SI{e-2}{\newton}$, vertical e para baixo.
Determine:
\begin{itensq}
\item a intensidade, direção e sentido do campo em P;
\item a força que atuaria sobre uma carga de $\SI{3}{\milli\coulomb}$ colocada em P.
\end{itensq}
\resposta{(a) $E=\pot{1.0}{7}\,\si{\newton\per\coulomb}$, vertical para baixo;
(b) $F = \pot{3}{4}\,\si{\newton}$}
\resolucao{(a) $E = \dfrac{F}{q} = \dfrac{\pot{1}{-2}}{\pot{1}{-9}}
= \pot{1.0}{7}\,\si{\newton\per\coulomb}$. Como a carga de prova é positiva, o
campo tem o \emph{mesmo} sentido da força: vertical, para baixo.\\
(b) $F = qE = \pot{3}{-3}\cdot\pot{1}{7} = \pot{3}{4}\,\si{\newton}$,
também vertical para baixo (carga positiva).}
\end{questao}

\begin{questao}[2]{}
Uma carga positiva $Q = \SI{4.0}{\micro\coulomb}$ está no ar. Determine a
intensidade, a direção e o sentido do campo elétrico em um ponto M, a
$\SI{20}{\centi\meter}$ da carga.
\dados{$\ke = \SI{9e9}{\newton\meter\squared\per\coulomb\squared}$.}
\resposta{$E = \pot{9.0}{5}\,\si{\newton\per\coulomb}$, radial, afastando-se de $Q$}
\resolucao{$E = \ke\dfrac{Q}{d^2}
= 9\cdot10^9\cdot\dfrac{\pot{4}{-6}}{(0{,}2)^2} = \pot{9.0}{5}\,\si{\newton\per\coulomb}$.
A direção é a da reta que une $Q$ a M, e o sentido é \emph{para fora} de $Q$
(carga positiva).}
\end{questao}

\begin{questao}[2]{}
Duas cargas puntiformes, $Q$ e $4Q$ (de mesmo sinal), estão fixas no vácuo. O
módulo do campo elétrico resultante é \textbf{nulo} em um ponto sobre a reta que
as une. Esse ponto está tal que sua distância à carga menor, $d_1$, e à carga
maior, $d_2$, satisfazem:
\begin{alternativas}[2]
\item $d_2 = d_1$
\item $d_2 = 2d_1$
\item $d_2 = 4d_1$
\item $d_1 = 2d_2$
\item $d_1 = 4d_2$
\end{alternativas}
\resposta{(b)}
\resolucao{O campo se anula \emph{entre} as cargas, onde os módulos se igualam:
\[
\ke\frac{Q}{d_1^2}=\ke\frac{4Q}{d_2^2}
\ \Rightarrow\
\frac{d_2^2}{d_1^2}=4
\ \Rightarrow\
d_2 = 2d_1.
\]
O ponto de campo nulo fica \emph{mais próximo} da carga menor.}
\end{questao}

\begin{questao}[2]{}
Duas cargas de mesmo módulo $Q$ estão fixas, separadas por $\SI{20}{\centi\meter}$.
Analise o campo elétrico resultante no ponto médio M do segmento que as une, em
dois casos:

\circuitoq{%
\begin{tikzpicture}[scale=0.9,line width=0.8pt,>=Stealth]
 % caso 1: +Q e -Q
 \begin{scope}
   \shade[ball color=Vcol!25] (-1.6,0) circle (0.25); \node at (-1.6,0){\footnotesize$+$};
   \shade[ball color=Icol!25] (1.6,0) circle (0.25);   \node at (1.6,0){\footnotesize$-$};
   \fill[laranja] (0,0) circle (0.06); \node[above,font=\footnotesize] at (0,0.12){M};
   \draw[Vcol,->,thick] (0,-0.5)--(1.1,-0.5); \node[Vcol,font=\footnotesize] at (0.55,-0.8){$\vec E_1$};
   \draw[Vcol,->,thick] (0,-1.1)--(1.1,-1.1); \node[Vcol,font=\footnotesize] at (0.55,-1.4){$\vec E_2$};
   \node[font=\footnotesize] at (0,1.0){Caso I: $+Q$ e $-Q$};
 \end{scope}
\end{tikzpicture}}{Campo no ponto médio entre duas cargas.}

No \textbf{Caso I} ($+Q$ e $-Q$) e no \textbf{Caso II} (ambas $+Q$), o campo
resultante em M é, respectivamente:
\begin{alternativas}
\item nulo em I e nulo em II.
\item não-nulo em I e nulo em II.
\item nulo em I e não-nulo em II.
\item não-nulo em ambos.
\item nulo em ambos.
\end{alternativas}
\resposta{(b)}
\resolucao{Em M, cada carga produz um campo de mesmo módulo (mesma distância,
mesmo $|Q|$).\\
\textbf{Caso I ($+Q$, $-Q$):} os dois campos apontam no \emph{mesmo} sentido (para
longe da positiva e em direção à negativa), então se \textbf{somam}: campo
resultante \emph{não-nulo}, $E = 2\cdot\dfrac{\ke Q}{d^2}$.\\
\textbf{Caso II ($+Q$, $+Q$):} os campos apontam em sentidos \emph{opostos} e se
\textbf{cancelam}: campo resultante \emph{nulo}.\\
Cuidado: para o \emph{potencial} a conclusão se inverte (nulo no caso I, não-nulo
no caso II), porque potencial é escalar e campo é vetorial.}
\end{questao}

\blocoaprofundamento

\begin{questao}[3]{}
As cargas $q_1 = \SI{20}{\coulomb}$ e $q_2 = \SI{64}{\coulomb}$ estão fixas no
vácuo nos pontos A e B, separados por $\SI{9}{\meter}$. Determine a posição, sobre
o segmento AB, onde o campo elétrico resultante é nulo.
\dados{$\ke = \SI{9e9}{\newton\meter\squared\per\coulomb\squared}$.}
\resposta{A $\SI{4}{\meter}$ de A (e $\SI{5}{\meter}$ de B)}
\resolucao{Como as cargas têm o mesmo sinal, o campo se anula entre elas. Seja $x$
a distância a partir de A; a distância a B é $9 - x$. Igualando os módulos:
\[
\ke\frac{q_1}{x^2}=\ke\frac{q_2}{(9-x)^2}
\ \Rightarrow\
\frac{20}{x^2}=\frac{64}{(9-x)^2}.
\]
Extraindo a raiz (grandezas positivas):
$\dfrac{9-x}{x}=\sqrt{\dfrac{64}{20}}=\dfrac{8}{\sqrt{20}}=\dfrac{8}{2\sqrt5}
=\dfrac{4}{\sqrt5}$. Resolvendo,
$9 - x = \dfrac{4}{\sqrt5}x \Rightarrow x = \dfrac{9}{1+4/\sqrt5}\approx\SI{4}{\meter}$.
O ponto neutro fica a cerca de \SI{4}{\meter} de A e \SI{5}{\meter} de B --- mais
perto da carga menor, como sempre ocorre.}
\end{questao}

\begin{questao}[3]{}
A intensidade do campo elétrico de uma carga puntiforme positiva, em função da
distância $d$, obedece ao gráfico $E \times d$ abaixo (curva do tipo $E = c/d^2$).
Sabendo que a $\SI{3}{\meter}$ o campo vale $\SI{80}{\newton\per\coulomb}$,
determine o campo a $\SI{6}{\meter}$.

\circuitoqq{%
\begin{tikzpicture}
 \begin{axis}[width=9cm,height=5cm,xlabel={$d$ (m)},ylabel={$E$ (N/C)},
   xmin=0,xmax=8,ymin=0,ymax=90,xtick={0,2,3,4,6,8},ytick={0,20,40,60,80},
   grid=both,axis lines=left,domain=1.6:8,samples=100]
   \addplot[Vcol,very thick] {720/(x^2)};
   \addplot[only marks,mark=*,Vcol] coordinates {(3,80)(6,20)};
   \node[Vcol,anchor=south west,font=\footnotesize] at (axis cs:3,80) {$(3,\,80)$};
   \node[Vcol,anchor=south west,font=\footnotesize] at (axis cs:6,20) {$(6,\,20)$};
 \end{axis}
\end{tikzpicture}}{Campo elétrico em função da distância ($E \propto 1/d^2$).}

\resposta{$E = \SI{20}{\newton\per\coulomb}$}
\resolucao{Como $E \propto \dfrac{1}{d^2}$, dobrar a distância (de \SI{3}{} para
\SI{6}{\meter}) divide o campo por 4:
\[
E(6)=\frac{E(3)}{4}=\frac{80}{4}=\SI{20}{\newton\per\coulomb}.
\]
Formalmente, $E(3)\cdot3^2 = E(6)\cdot6^2$, ou seja,
$80\cdot9 = E(6)\cdot36 \Rightarrow E(6) = \SI{20}{\newton\per\coulomb}$.}
\end{questao}

\begin{questao}[3]{}
Duas cargas $q_1 = \SI{+9}{\micro\coulomb}$ e $q_2 = \SI{+4}{\micro\coulomb}$
estão fixas e separadas por $\SI{10}{\centi\meter}$. Determine a que distância da
carga $q_1$, sobre a reta que as une, o campo elétrico resultante é nulo.
\resposta{A $\SI{6}{\centi\meter}$ de $q_1$}
\resolucao{Cargas de mesmo sinal: o campo se anula \emph{entre} elas. Seja $x$ a
distância a $q_1$; a distância a $q_2$ é $(10 - x)$. Igualando os módulos dos
campos:
\[
\frac{\ke q_1}{x^2}=\frac{\ke q_2}{(10-x)^2}
\ \Rightarrow\
\frac{9}{x^2}=\frac{4}{(10-x)^2}.
\]
Extraindo a raiz: $\dfrac{3}{x}=\dfrac{2}{10-x} \Rightarrow 3(10-x)=2x
\Rightarrow 30 = 5x \Rightarrow x = \SI{6}{\centi\meter}$.
O ponto neutro fica mais perto da carga \emph{menor} ($q_2$), a
\SI{4}{\centi\meter} dela.}
\end{questao}

\begin{questao}[3]{}
Duas cargas puntiformes, $+4Q$ e $-Q$, estão fixas no vácuo, separadas por uma
distância $d$.
\begin{itensq}
\item Mostre que existe um ponto sobre a reta que as une onde o campo elétrico
      resultante é \textbf{nulo}, e determine sua posição.
\item Explique por que esse ponto está \emph{fora} do segmento que une as cargas,
      e do lado da carga de \emph{menor} módulo.
\item Nesse mesmo ponto, o potencial elétrico é nulo? Justifique.
\end{itensq}
\resposta{(a) a $2d$ da carga $+4Q$ (isto é, a $d$ além de $-Q$);
(b) e (c) ver comentário}
\resolucao{\textbf{(a)} Seja $x$ a distância medida a partir de $+4Q$, com o ponto
além de $-Q$ (logo a distância a $-Q$ é $x-d$). Igualando os módulos:
\[
\frac{\ke\,4Q}{x^{2}}=\frac{\ke\,Q}{(x-d)^{2}}
\;\Longrightarrow\;
4(x-d)^2 = x^2
\;\Longrightarrow\;
2(x-d)=\pm x.
\]
A raiz fisicamente válida é $2x-2d = x \Rightarrow \boxed{x = 2d}$, ou seja, a
$2d$ da carga maior e a $d$ além da carga menor.
(A outra raiz, $x = \tfrac{2d}{3}$, cai \emph{entre} as cargas, onde os campos
apontam no \emph{mesmo} sentido e portanto não podem se cancelar --- deve ser
descartada.)

\textbf{(b)} Entre cargas de sinais \emph{opostos}, os dois vetores campo apontam
no mesmo sentido (saindo da positiva, entrando na negativa) e se \emph{somam};
nunca se anulam ali. Fora do segmento, eles têm sentidos contrários, e o
cancelamento só é possível onde o campo da carga \emph{menor} --- que está mais
próxima --- consegue igualar o da maior. Por isso o ponto neutro fica sempre do
lado da carga de menor módulo.

\textbf{(c) Não.} Potencial é grandeza \emph{escalar}:
\[
V = \ke\frac{4Q}{2d}+\ke\frac{(-Q)}{d}
= \frac{2\ke Q}{d}-\frac{\ke Q}{d}
= \frac{\ke Q}{d}\neq 0.
\]
\textbf{Conclusão conceitual:} campo nulo e potencial nulo são condições
\emph{independentes}. Neste caso, $\vec E = 0$ mas $V \neq 0$. No ponto médio de um
dipolo simétrico ocorre o oposto: $V = 0$ mas $\vec E \neq 0$.}
\end{questao}

\begin{questao}[3]{}
Um elétron entra com velocidade horizontal $v_0=\SI{2e7}{\meter\per\second}$ na
região entre duas placas paralelas de comprimento $L=\SI{5}{\centi\meter}$, onde
existe campo elétrico uniforme $E=\SI{e4}{\volt\per\meter}$, perpendicular à
velocidade.

\circuitoq{%
\begin{tikzpicture}[scale=1,line width=0.8pt,>=Stealth]
 \draw[cinza!70,line width=2.2pt] (0,1.3)--(4.2,1.3);
 \draw[cinza!70,line width=2.2pt] (0,-1.3)--(4.2,-1.3);
 \node[above,font=\footnotesize] at (2.1,1.35) {$+$\ placa};
 \node[below,font=\footnotesize] at (2.1,-1.35) {$-$\ placa};
 \foreach \x in {0.6,1.5,2.4,3.3}
   \draw[->,Vcol!70] (\x,1.2)--(\x,-1.2);
 \node[Vcol,font=\footnotesize] at (3.85,0.55) {$\vec E$};
 \draw[->,laranja,line width=1.2pt] (-1.4,0)--(0,0)
       node[midway,above,font=\footnotesize]{$v_0$};
 \draw[laranja,line width=1.2pt,domain=0:4.2,samples=40,smooth]
       plot (\x,{0.055*\x*\x});
 \draw[->,laranja,line width=1.2pt] (4.2,{0.055*4.2*4.2})--(5.3,{0.055*4.2*4.2+0.55});
 \draw[<->,cinza] (0,-1.75)--(4.2,-1.75) node[midway,below,font=\footnotesize]{$L$};
 \draw[<->,cinza] (4.5,0)--(4.5,{0.055*4.2*4.2}) node[midway,right,font=\footnotesize]{$y$};
\end{tikzpicture}}{Deflexão de um elétron em campo elétrico uniforme.}

\begin{itensq}
\item Calcule a aceleração do elétron.
\item Determine o tempo de permanência entre as placas e o desvio vertical $y$.
\item Calcule o ângulo de saída em relação à direção inicial.
\end{itensq}
\dados{$e=\SI{1.6e-19}{\coulomb}$; $m_e=\SI{9.11e-31}{\kilogram}$.}
\resposta{(a) $a\approx\SI{1.76e15}{\meter\per\second\squared}$;
(b) $t=\SI{2.5}{\nano\second}$, $y\approx\SI{5.5}{\milli\meter}$;
(c) $\theta\approx\SI{12.4}{\degree}$}
\resolucao{Este é um \textbf{lançamento de projétil} eletrostático: movimento
uniforme na horizontal e uniformemente acelerado na vertical.\\
(a) $a=\dfrac{eE}{m_e}=\dfrac{(\pot{1.6}{-19})(\pot{1}{4})}{\pot{9.11}{-31}}
\approx\pot{1.76}{15}\,\si{\meter\per\second\squared}$
--- cerca de $10^{14}$ vezes a gravidade, o que justifica desprezar o peso.\\
(b) Na horizontal (MU):
$t=\dfrac{L}{v_0}=\dfrac{0{,}05}{\pot{2}{7}}=\pot{2.5}{-9}\,\si{\second}
=\SI{2.5}{\nano\second}$.\\
Na vertical (MUV, partindo do repouso):
\[
y=\tfrac12 a t^{2}=\tfrac12(\pot{1.76}{15})(\pot{2.5}{-9})^{2}
\approx\pot{5.5}{-3}\,\si{\meter}=\SI{5.5}{\milli\meter}.
\]
(c) A velocidade vertical na saída é $v_y = a\,t \approx
\pot{4.4}{6}\,\si{\meter\per\second}$, logo
\[
\tan\theta=\frac{v_y}{v_0}=\frac{\pot{4.4}{6}}{\pot{2}{7}}=0{,}22
\;\Longrightarrow\;
\theta\approx\SI{12.4}{\degree}.
\]
\textbf{Aplicação:} é exatamente esse o mecanismo de deflexão dos antigos tubos de
raios catódicos (TV e osciloscópios) e das impressoras a jato de tinta, que
defletem gotículas carregadas para formar caracteres.}
\end{questao}

\begin{questao}[3]{}
Quatro cargas de módulo $q=\SI{1}{\micro\coulomb}$ ocupam os vértices de um
quadrado de lado $a=\SI{10}{\centi\meter}$. Compare o \textbf{campo elétrico no
centro} em três configurações:
\begin{itensq}
\item todas positivas;
\item duas positivas e duas negativas, alternadas ($+,-,+,-$);
\item duas positivas e duas negativas, adjacentes ($+,+,-,-$).
\end{itensq}
\dados{$\ke=\SI{9e9}{}$; $\sqrt2\approx1{,}41$.}
\resposta{(a) $\vec E=0$; (b) $\vec E=0$; (c) $E\approx\pot{5.1}{6}\,\si{\newton\per\coulomb}$}
\resolucao{Cada carga dista $\dfrac{a\sqrt2}{2}=\SI{0.0707}{\meter}$ do centro,
produzindo individualmente
\[
E_1=\frac{\ke q}{(a\sqrt2/2)^{2}}=\frac{9\cdot10^{9}\cdot\pot{1}{-6}}{0{,}005}
=\pot{1.8}{6}\,\si{\newton\per\coulomb}.
\]
\textbf{(a) Todas positivas.} Os quatro vetores apontam radialmente para fora,
aos pares diametralmente opostos: cancelam-se dois a dois. $\vec E = 0$.\\
\textbf{(b) Alternadas ($+,-,+,-$).} As cargas diametralmente opostas têm o
\emph{mesmo} sinal, e seus campos continuam se cancelando aos pares.
$\vec E = 0$ novamente.\\
\textbf{(c) Adjacentes ($+,+,-,-$).} Agora cada par diametralmente oposto tem
sinais \emph{contrários}: os campos se \textbf{somam} em vez de cancelar. Os dois
pares contribuem na mesma direção (perpendicular à linha que separa as positivas
das negativas):
\[
E=2\,(E_1\sqrt2)=2\sqrt2\,E_1\approx 2{,}83\cdot\pot{1.8}{6}
\approx\pot{5.1}{6}\,\si{\newton\per\coulomb}.
\]
\textbf{Lição central.} Em (a) e (b), configurações \emph{muito} diferentes
produzem o mesmo resultado nulo. Em (c), com as mesmas quatro cargas, apenas
reposicionadas, o campo é máximo. \emph{O campo resultante depende da geometria
tanto quanto dos valores das cargas} --- e a simetria é a ferramenta que revela
isso sem nenhum cálculo pesado.}
\end{questao}

\begin{questao}[3]{}
Uma carga $-Q$ está na origem e uma carga $+4Q$ está em $x=d$.
\begin{itensq}
\item Determine o ponto sobre o eixo $x$ onde o campo resultante é nulo.
\item Explique por que ele fica \emph{fora} do segmento e do lado da carga de
      menor módulo.
\item Nesse ponto, quanto vale o potencial elétrico?
\end{itensq}
\resposta{(a) $x=-d$ (a uma distância $d$ à esquerda de $-Q$);
(b) e (c) ver comentário; $V=\dfrac{\ke Q}{d}$}
\resolucao{(a) Como as cargas têm sinais \emph{opostos}, entre elas os campos
apontam no mesmo sentido e se somam --- o cancelamento só pode ocorrer fora do
segmento. Seja $x<0$ a posição procurada; a distância a $-Q$ é $|x|$ e a $+4Q$ é
$|x|+d$:
\[
\frac{\ke Q}{|x|^{2}}=\frac{\ke\,4Q}{(|x|+d)^{2}}
\;\Longrightarrow\;
\frac{|x|+d}{|x|}=2
\;\Longrightarrow\;
|x|=d.
\]
O ponto está em $x=-d$.\\
(b) Do lado da carga \emph{maior}, seu campo sempre domina (está mais perto
\emph{e} é mais intensa) --- nunca poderia ser equilibrado. Do lado da carga
menor, a proximidade compensa a menor intensidade, e existe exatamente um ponto
de equilíbrio.\\
(c) O potencial é escalar:
\[
V=\ke\frac{(-Q)}{d}+\ke\frac{4Q}{2d}
=-\frac{\ke Q}{d}+\frac{2\ke Q}{d}
=\frac{\ke Q}{d}\neq0.
\]
\textbf{Mais um exemplo de que $\vec E=0$ não implica $V=0$} --- e vice-versa. O
campo mede a \emph{taxa de variação} do potencial, não o seu valor.}
\end{questao}

\begin{questao}[3]{}
Um anel isolante de raio $R$ está uniformemente carregado com carga total $Q$.
Considere um ponto P sobre o eixo do anel, a uma distância $x$ do centro.
\begin{itensq}
\item Explique, por simetria, por que o campo em P é paralelo ao eixo.
\item Sabendo que $E(x)=\dfrac{\ke\,Q\,x}{\left(x^{2}+R^{2}\right)^{3/2}}$,
      determine $E$ no centro do anel e no limite $x\gg R$.
\item Mostre que o campo é \emph{máximo} em $x = \dfrac{R}{\sqrt2}$.
\end{itensq}
\resposta{(b) $E(0)=0$ e $E\to\ke Q/x^{2}$; (c) demonstração}
\resolucao{(a) \textbf{Argumento de simetria.} Para cada elemento de carga do
anel existe outro diametralmente oposto. As componentes de campo
\emph{perpendiculares} ao eixo desses dois elementos são iguais e opostas,
cancelando-se. Restam apenas as componentes \emph{paralelas} ao eixo, que se
somam. Logo $\vec E$ é axial.

(b) \textbf{No centro} ($x=0$): substituindo,
\[
E(0)=\frac{\ke Q\cdot 0}{R^{3}}=0.
\]
Faz sentido: no centro, as contribuições de elementos opostos se cancelam
\emph{integralmente}.\\
\textbf{Longe} ($x\gg R$): então $\left(x^{2}+R^{2}\right)^{3/2}\approx x^{3}$ e
\[
E\approx\frac{\ke Q x}{x^{3}}=\frac{\ke Q}{x^{2}},
\]
o campo de uma carga puntiforme --- como esperado, pois de longe o anel "vira"
um ponto.

(c) \textbf{Máximo.} Derivando $E(x)$ e igualando a zero:
\[
\frac{\mathrm{d}E}{\mathrm{d}x}
=\ke Q\,\frac{\left(x^{2}+R^{2}\right)^{3/2}-x\cdot\tfrac32\left(x^{2}+R^{2}\right)^{1/2}(2x)}
{\left(x^{2}+R^{2}\right)^{3}}=0.
\]
O numerador, dividido por $\left(x^{2}+R^{2}\right)^{1/2}$, dá
\[
\left(x^{2}+R^{2}\right)-3x^{2}=0
\;\Longrightarrow\;
R^{2}=2x^{2}
\;\Longrightarrow\;
\boxed{\,x=\frac{R}{\sqrt2}\,}
\]
\textbf{Interpretação física do máximo.} Perto do centro, as contribuições
axiais são pequenas (os vetores são quase perpendiculares ao eixo). Muito longe,
a distância cresce e o campo decai. Entre os dois regimes existe um ponto ótimo,
a $x=R/\sqrt2\approx0{,}71R$.\\
\textbf{Onde isso é usado:} essa geometria é a base das \emph{bobinas de
Helmholtz} (versão magnética do mesmo cálculo) e dos anéis focalizadores em
aceleradores de partículas e microscópios eletrônicos.}
\end{questao}

\blocodesafio

\begin{questao}[4]{}
Uma haste isolante fina de comprimento $L$ está uniformemente carregada com carga
total $Q$ (densidade linear $\lambda=Q/L$). Considere um ponto P sobre o
prolongamento da haste, a uma distância $a$ da extremidade mais próxima.
\begin{itensq}
\item Monte a integral que fornece o campo elétrico em P.
\item Mostre que $E=\dfrac{\ke\,Q}{a\,(a+L)}$.
\item Verifique o comportamento no limite $a \gg L$ e interprete.
\end{itensq}
\resposta{(b) demonstração; (c) $E\to \ke Q/a^{2}$ (carga puntiforme)}
\resolucao{(a) Tome um elemento de comprimento $\mathrm{d}x$ à distância $x$ de P
(com $x$ variando de $a$ até $a+L$). Sua carga é
$\mathrm{d}q=\lambda\,\mathrm{d}x$ e contribui com
\[
\mathrm{d}E=\frac{\ke\,\mathrm{d}q}{x^{2}}=\frac{\ke\lambda\,\mathrm{d}x}{x^{2}}.
\]
Como todos os elementos estão alinhados com P, as contribuições têm a
\emph{mesma direção} e somam-se escalarmente:
\[
E=\int_{a}^{a+L}\frac{\ke\lambda}{x^{2}}\,\mathrm{d}x.
\]
(b) Integrando:
\[
E=\ke\lambda\left[-\frac{1}{x}\right]_{a}^{a+L}
=\ke\lambda\left(\frac{1}{a}-\frac{1}{a+L}\right)
=\ke\lambda\,\frac{L}{a(a+L)}.
\]
Como $\lambda L = Q$:
\[
\boxed{\;E=\frac{\ke\,Q}{a\,(a+L)}\;}
\]
(c) \textbf{Limite $a\gg L$:} então $a+L\approx a$ e
\[
E\approx\frac{\ke Q}{a^{2}},
\]
que é exatamente o campo de uma \textbf{carga puntiforme} $Q$.\\
\textbf{Interpretação:} de longe, qualquer distribuição finita de carga
"parece" puntiforme. Esse teste de consistência --- verificar se a expressão
obtida reproduz o caso conhecido em algum limite --- é uma das melhores formas de
validar um resultado. \emph{Se o limite não fechasse, haveria erro na
integração.}\\
\textbf{Limite oposto ($a\ll L$):} $E\approx\dfrac{\ke Q}{aL}=\dfrac{\ke\lambda}{a}$
--- campo de um fio "infinito" visto de perto, decaindo com $1/a$ em vez de
$1/a^{2}$.}
\end{questao}

\begin{questao}[4]{}
Um dipolo elétrico é formado por cargas $+q$ e $-q$ separadas por uma pequena
distância $2\ell$. Define-se o \textbf{momento de dipolo} $p=q\,(2\ell)$.
\begin{itensq}
\item Mostre que, sobre a mediatriz do dipolo, a uma distância $r\gg\ell$ do
      centro, o campo vale $E\approx\dfrac{\ke\,p}{r^{3}}$.
\item Por que o campo do dipolo decai mais rápido que o de uma carga isolada?
\item Um dipolo colocado em campo uniforme sofre força resultante? E torque?
\end{itensq}
\resposta{(a) demonstração; (b) e (c) ver comentário}
\resolucao{(a) Sobre a mediatriz, cada carga dista
$d=\sqrt{r^{2}+\ell^{2}}$ do ponto e produz campo de módulo
$E_1=\dfrac{\ke q}{r^{2}+\ell^{2}}$. As componentes \emph{paralelas} à mediatriz
se cancelam (simetria); restam as componentes ao longo do eixo do dipolo, cada uma
valendo $E_1\cos\alpha$, com $\cos\alpha=\dfrac{\ell}{\sqrt{r^{2}+\ell^{2}}}$:
\[
E=2E_1\cos\alpha
=\frac{2\ke q\,\ell}{\left(r^{2}+\ell^{2}\right)^{3/2}}.
\]
Para $r\gg\ell$, $\left(r^{2}+\ell^{2}\right)^{3/2}\approx r^{3}$, e como
$p=2q\ell$:
\[
\boxed{\;E\approx\frac{\ke\,p}{r^{3}}\;}
\]
(b) Porque a carga \emph{total} do dipolo é nula. De longe, os campos das duas
cargas quase se cancelam --- sobra apenas o efeito da pequena separação entre
elas, que é de ordem superior. Daí o decaimento $1/r^{3}$ em vez de $1/r^{2}$.\\
(c) Em campo \textbf{uniforme}:
\begin{itemize}[leftmargin=1.6em,itemsep=1pt,topsep=2pt]
\item \textbf{Força resultante nula}: as forças $+qE$ e $-qE$ sobre as duas cargas
      têm mesmo módulo e sentidos opostos.
\item \textbf{Torque não nulo}: como elas atuam em \emph{pontos diferentes},
      formam um binário de momento
      \[
      \tau = p\,E\sin\theta,
      \]
      que tende a \emph{alinhar} o dipolo com o campo.
\end{itemize}
\textbf{Aplicações:} é esse torque que alinha moléculas polares (como a água) em
um campo --- princípio do forno de micro-ondas, onde o campo oscilante faz as
moléculas girarem e o atrito molecular gera calor. Em campo \emph{não} uniforme,
surge também força resultante --- razão pela qual um objeto neutro é atraído por
um corpo eletrizado.}
\end{questao}

% BEGIN BANCO COMPLEMENTAR Campo elétrico
% =====================================================================
%  Banco complementar --- Campo Eletrico  (REVISADO)
%  Substitui a versao gerada por template (8 clones em N1, 11 em N2,
%  7 em N3 e 8 questoes conceituais marcadas como N4).
%  O cap01b e denso e ja cobre: E de carga puntiforme, proporcionalidade
%  com a distancia, carga de prova com forca dada, direcao e sentido do
%  campo, PONTO DE CAMPO NULO (tres vezes: Q e 4Q; 9 e 4 uC; -Q e +4Q),
%  duas cargas de mesmo modulo a 20 cm, grafico E x d, ELETRON ENTRE
%  PLACAS PARALELAS com deflexao, quadrado de quatro cargas, anel
%  isolante no eixo, haste finita por integracao (N4) e momento dipolar
%  (N4).
%  Este banco evita esses casos e explora: E=0 com V nao nulo, linhas de
%  campo e por que nao se cruzam, PROJETO de placas paralelas, cavidade
%  em condutor, tracado numerico de linhas, esfera isolante macica por
%  Gauss, par de placas com +-sigma e o limite imposto pela rigidez
%  dieletrica.
%  Distribuicao mantida: 8 N1 + 11 N2 + 7 N3 + 8 N4 = 34
% =====================================================================

\subsubsection*{Banco complementar --- Campo elétrico}

\begin{atencao}[Organização do banco]
O campo é \textbf{vetorial}: contribuições somam-se por componentes, e podem
cancelar-se. Ele é definido por $\vec E=\vec F/q_0$, mas \emph{não depende} da
carga de prova --- é propriedade do ponto. Duas expressões dominam o tópico:
$E=\ke Q/r^{2}$ para carga puntiforme e $E=\sigma/\varepsilon_0$ para plano
carregado, esta última \emph{independente da distância}. O N4 trata de projeto,
condutores, distribuições contínuas e limites físicos.
\end{atencao}

\blocofixacao

\begin{questao}[1]{}
Determine o módulo do campo elétrico produzido por $Q=+\SI{3}{\micro\coulomb}$
a $\SI{30}{\centi\meter}$ de distância.
\dados{$\ke=\SI{9e9}{\newton\meter\squared\per\coulomb\squared}$}
\resposta{$E=\SI{300}{\kilo\newton\per\coulomb}$}
\resolucao{
\[
E=\frac{\ke Q}{r^{2}}=\frac{\pot{9}{9}\times\pot{3}{-6}}{(0{,}30)^{2}}
=\frac{\pot{2.7}{4}}{0{,}09}=\pot{3}{5}\,\si{\newton\per\coulomb}.
\]
O campo aponta \textbf{para longe} da carga, por ela ser positiva.\\
\textbf{Atenção à potência:} o campo cai com $1/r^{2}$, o potencial com
$1/r$. Dobrar a distância divide $E$ por quatro e $V$ por dois.}
\end{questao}

\begin{questao}[1]{}
Determine o campo produzido por $Q=-\SI{5}{\micro\coulomb}$ a
$\SI{50}{\centi\meter}$, indicando o sentido.
\resposta{$E=\SI{180}{\kilo\newton\per\coulomb}$, apontando \emph{para} a carga}
\resolucao{
\[
E=\frac{\pot{9}{9}\times\pot{5}{-6}}{0{,}25}
=\pot{1.8}{5}\,\si{\newton\per\coulomb}.
\]
\textbf{Sentido:} cargas negativas \emph{atraem} cargas de prova positivas, logo
o campo aponta \textbf{para} a carga fonte.\\
\textbf{Regra mnemônica:} o campo \emph{sai} do positivo e \emph{entra} no
negativo --- é a mesma regra que orienta as linhas de campo.}
\end{questao}

\begin{questao}[1]{}
Uma carga $q=+\SI{2}{\micro\coulomb}$ é colocada em um ponto onde
$E=\SI{5000}{\newton\per\coulomb}$. Determine a força sobre ela.
\resposta{$F=\SI{10}{\milli\newton}$, no sentido do campo}
\resolucao{
\[
F=qE=\pot{2}{-6}\times5000=\pot{1}{-2}\,\si{\newton}=\SI{10}{\milli\newton}.
\]
Como a carga é positiva, a força tem o \emph{mesmo} sentido de $\vec E$. Se
fosse negativa, teria sentido oposto --- e apenas o sentido mudaria, não o
módulo.}
\end{questao}

\begin{questao}[1]{}
Determine a carga que produz $E=\SI{2000}{\newton\per\coulomb}$ a
$\SI{30}{\centi\meter}$.
\resposta{$Q=\SI{20}{\nano\coulomb}$}
\resolucao{
\[
Q=\frac{Er^{2}}{\ke}=\frac{2000\times0{,}09}{\pot{9}{9}}
=\pot{2}{-8}\,\si{\coulomb}=\SI{20}{\nano\coulomb}.
\]
Uma carga de nanocoulombs --- ordem de grandeza típica da eletrização por
atrito --- já produz milhares de newtons por coulomb a trinta centímetros.}
\end{questao}

\begin{questao}[1]{}
A que distância de $Q=+\SI{8}{\micro\coulomb}$ o campo vale
$\SI{200}{\kilo\newton\per\coulomb}$?
\resposta{$r=\SI{0.6}{\meter}$}
\resolucao{
\[
r=\sqrt{\frac{\ke Q}{E}}=\sqrt{\frac{\pot{9}{9}\times\pot{8}{-6}}{\pot{2}{5}}}
=\sqrt{0{,}36}=\SI{0.6}{\meter}.
\]
\textbf{Superfície de campo constante:} todos os pontos a
$\SI{0.6}{\meter}$ da carga têm o mesmo \emph{módulo} de campo --- mas
\emph{direções} diferentes, pois o campo é radial. Diferentemente das
equipotenciais, essas superfícies não têm significado especial.}
\end{questao}

\begin{questao}[1]{}
O campo elétrico em um ponto, medido com carga de prova $q_0$, vale $E$. Se a
carga de prova for \textbf{dobrada}, o valor medido:
\begin{alternativas}[2]
\item dobra
\item cai à metade
\item não se altera
\item depende do sinal de $q_0$
\end{alternativas}
\resposta{(c)}
\resolucao{Por definição, $\vec E=\vec F/q_0$. Dobrar $q_0$ dobra a força
\emph{e} o denominador --- a razão permanece.\\
\textbf{Ponto conceitual:} o campo é uma propriedade \textbf{do ponto do
espaço}, criada pelas cargas-fonte. A carga de prova apenas o \emph{revela};
ela não o cria nem o modifica.\\
\textbf{Ressalva real:} se $q_0$ for grande demais, ela perturba as
cargas-fonte (por indução, em condutores) e a medida deixa de valer. Daí a
definição rigorosa usar o limite $q_0\to0$.}
\end{questao}

\begin{questao}[1]{}
Mostre que $\si{\newton\per\coulomb}$ e $\si{\volt\per\meter}$ são a mesma
unidade.
\resposta{$\si{\volt}=\si{\joule\per\coulomb}$ e
$\si{\joule}=\si{\newton\meter}$}
\resolucao{
\[
\frac{\si{\volt}}{\si{\meter}}
=\frac{\si{\joule}}{\si{\coulomb}\cdot\si{\meter}}
=\frac{\si{\newton}\cdot\si{\meter}}{\si{\coulomb}\cdot\si{\meter}}
=\frac{\si{\newton}}{\si{\coulomb}}.\ \checkmark
\]
\textbf{As duas leituras do campo.} $\si{\newton\per\coulomb}$ enfatiza a
\emph{força} por unidade de carga; $\si{\volt\per\meter}$ enfatiza a
\emph{taxa de queda do potencial} por unidade de comprimento. São descrições
equivalentes de $\vec E=-\nabla V$.\\
Na prática, usa-se $\si{\volt\per\meter}$ quando há tensões envolvidas (placas,
cabos) e $\si{\newton\per\coulomb}$ em problemas de dinâmica de partículas.}
\end{questao}

\begin{questao}[1]{}
Uma carga negativa é colocada em uma região de campo elétrico uniforme. Sobre
seu movimento inicial, é correto afirmar que ela:
\begin{alternativas}[2]
\item acelera no sentido do campo
\item acelera no sentido oposto ao campo
\item permanece em repouso
\item move-se perpendicularmente ao campo
\end{alternativas}
\resposta{(b)}
\resolucao{A força vale $\vec F=q\vec E$ com $q<0$, logo $\vec F$ é
\textbf{antiparalela} a $\vec E$.\\
\textbf{Consequência energética:} a carga negativa move-se espontaneamente no
sentido de $V$ \emph{crescente} --- o oposto de uma carga positiva. Em ambos os
casos, a energia potencial $U=qV$ \emph{diminui}, como deve ocorrer em qualquer
movimento espontâneo.}
\end{questao}

\blocoaplicacao

\begin{questao}[2]{}
Uma carga $q_0=-\SI{1}{\micro\coulomb}$ é colocada em um campo uniforme de
$E=\SI{2000}{\newton\per\coulomb}$.
\begin{itensq}
\item Determine a força.
\item Determine sua aceleração, se a massa for $\SI{4}{\gram}$.
\end{itensq}
\resposta{$F=\SI{2}{\milli\newton}$, oposta a $\vec E$;
$a=\SI{0.5}{\meter\per\second\squared}$}
\resolucao{(a) $F=|q|E=\pot{1}{-6}(2000)=\pot{2}{-3}\,\si{\newton}$, com sentido
\emph{oposto} ao campo.\\
(b)
\[
a=\frac{F}{m}=\frac{\pot{2}{-3}}{\pot{4}{-3}}=\SI{0.5}{\meter\per\second\squared}.
\]
\textbf{Comparação com a gravidade:} $a$ é apenas $5\%$ de $g$. Para objetos
macroscópicos, a força elétrica só compete com o peso se a carga for
relativamente grande ou a massa muito pequena --- daí os experimentos
eletrostáticos clássicos usarem pedaços de papel e gotas de óleo.}
\end{questao}

\begin{questao}[2]{}
Um próton é abandonado em um campo uniforme de
$E=\SI{1e4}{\newton\per\coulomb}$. Determine sua aceleração.
\dados{$e=\pot{1.6}{-19}\,\si{\coulomb}$; $m_p=\pot{1.67}{-27}\,\si{\kilogram}$}
\resposta{$a=\pot{9.6}{11}\,\si{\meter\per\second\squared}$}
\resolucao{
\[
a=\frac{eE}{m_p}=\frac{(\pot{1.6}{-19})(\pot{1}{4})}{\pot{1.67}{-27}}
=\pot{9.58}{11}\,\si{\meter\per\second\squared}.
\]
\textbf{Cerca de $10^{11}$ vezes $g$} --- a gravidade é completamente
desprezível na dinâmica de partículas carregadas.\\
\textbf{Para o elétron}, a aceleração seria $1836$ vezes maior
($\pot{1.76}{15}\,\si{\meter\per\second\squared}$), pela razão das massas. É por
isso que, em tubos e aceleradores, são os elétrons que respondem primeiro.}
\end{questao}

\begin{questao}[2]{}
Duas cargas estão sobre o eixo $x$: $+\SI{4}{\micro\coulomb}$ em $x=0$ e
$-\SI{2}{\micro\coulomb}$ em $x=\SI{0.30}{\meter}$. Determine o campo em
$x=\SI{0.60}{\meter}$.
\resposta{$E=\SI{100}{\kilo\newton\per\coulomb}$, no sentido $-x$}
\resolucao{
\[
E_1=\frac{\pot{9}{9}(\pot{4}{-6})}{(0{,}60)^{2}}
=\pot{1}{5}\,\si{\newton\per\coulomb}\ (\text{sentido }+x);
\]
\[
E_2=\frac{\pot{9}{9}(\pot{2}{-6})}{(0{,}30)^{2}}
=\pot{2}{5}\,\si{\newton\per\coulomb}\ (\text{sentido }-x).
\]
\[
E=2\times10^{5}-1\times10^{5}=\pot{1}{5}\,\si{\newton\per\coulomb}\ (-x).
\]
\textbf{Observe:} a carga \emph{menor} produziu o campo \emph{maior}, por estar
duas vezes mais próxima --- e $1/r^{2}$ vence o fator $2$ de carga.
Compare os fatores: carga $\times\frac12$, distância$^{2}$
$\times\frac14$.}
\end{questao}

\begin{questao}[2]{}
Duas cargas iguais $Q=+\SI{2}{\micro\coulomb}$ estão em
$(\pm\SI{0.30}{\meter},\,0)$. Determine o campo no ponto
$(0,\ \SI{0.40}{\meter})$.
\resposta{$E=\SI{115.2}{\kilo\newton\per\coulomb}$, ao longo de $+y$}
\resolucao{Cada carga dista $r=\sqrt{0{,}30^{2}+0{,}40^{2}}=\SI{0.5}{\meter}$
(triângulo $3$--$4$--$5$):
\[
E_1=E_2=\frac{\pot{9}{9}(\pot{2}{-6})}{0{,}25}
=\pot{7.2}{4}\,\si{\newton\per\coulomb}.
\]
\textbf{Por simetria}, as componentes horizontais se cancelam e as verticais se
somam, cada uma multiplicada por $\cos\theta=0{,}40/0{,}50=0{,}8$:
\[
E=2(\pot{7.2}{4})(0{,}8)=\pot{1.152}{5}\,\si{\newton\per\coulomb}.
\]
\textbf{Uso da simetria:} reconhecê-la elimina metade do trabalho. Sempre que a
configuração for simétrica em relação ao ponto de interesse, identifique quais
componentes se anulam \emph{antes} de calcular.}
\end{questao}

\begin{questao}[2]{}
Duas placas paralelas separadas por $\SI{2}{\milli\meter}$ estão submetidas a
$\SI{600}{\volt}$. Determine o campo entre elas.
\resposta{$E=\SI{300}{\kilo\volt\per\meter}$}
\resolucao{
\[
E=\frac{U}{d}=\frac{600}{\pot{2}{-3}}=\pot{3}{5}\,\si{\volt\per\meter}.
\]
\textbf{Campo uniforme:} entre placas extensas e próximas, $E$ é praticamente o
mesmo em todos os pontos --- e independe da posição. É a configuração de
referência para acelerar e defletir partículas.\\
\textbf{Verificação de segurança:} $\SI{300}{\kilo\volt\per\meter}$ é um décimo
da rigidez do ar ($\SI{3}{\mega\volt\per\meter}$) --- há folga confortável.
Reduzir a separação para $\SI{0.2}{\milli\meter}$ com a mesma tensão provocaria
ruptura.}
\end{questao}

\begin{questao}[2]{}
Uma placa condutora extensa tem densidade superficial de carga
$\sigma=\SI{2}{\micro\coulomb\per\meter\squared}$. Determine o campo próximo à
sua superfície.
\dados{$\varepsilon_0=\pot{8.85}{-12}\,\si{\farad\per\meter}$}
\resposta{$E=\SI{226}{\kilo\newton\per\coulomb}$}
\resolucao{Para um condutor em equilíbrio,
\[
E=\frac{\sigma}{\varepsilon_0}=\frac{\pot{2}{-6}}{\pot{8.85}{-12}}
=\pot{2.26}{5}\,\si{\newton\per\coulomb}.
\]
\textbf{Fato notável:} o campo \textbf{não depende da distância} à placa (desde
que ela seja extensa e o ponto próximo). Diferentemente da carga puntiforme, não
há queda com $1/r^{2}$.\\
\textbf{Cuidado com a fórmula:} para um \emph{plano isolante} com carga nas duas
faces, o campo vale $\sigma/2\varepsilon_0$ --- metade. A diferença é que no
condutor toda a carga fica em uma face, e o campo interno é nulo.}
\end{questao}

\begin{questao}[2]{}
Compare o campo produzido por $Q_1=+\SI{9}{\micro\coulomb}$ a
$\SI{30}{\centi\meter}$ com o de $Q_2=+\SI{4}{\micro\coulomb}$ a
$\SI{20}{\centi\meter}$.
\resposta{$E_1=E_2=\SI{900}{\kilo\newton\per\coulomb}$ --- iguais}
\resolucao{
\[
E_1=\frac{\pot{9}{9}(\pot{9}{-6})}{0{,}09}=\pot{9}{5};
\qquad
E_2=\frac{\pot{9}{9}(\pot{4}{-6})}{0{,}04}=\pot{9}{5}\,\si{\newton\per\coulomb}.
\]
\textbf{Coincidência não é acaso:} as razões $Q/r^{2}$ são iguais, pois
$9/9=4/4$. Sempre que $Q\propto r^{2}$, o campo se mantém.\\
\textbf{Consequência:} se essas duas cargas estivessem em lados opostos de um
ponto, o campo resultante ali seria \emph{nulo} --- e essa é exatamente a
condição do "ponto de campo nulo" entre cargas de mesmo sinal.}
\end{questao}

\begin{questao}[2]{}
Uma carga em um ponto sofre campos de $\SI{3}{\kilo\newton\per\coulomb}$ na
direção $x$ e $\SI{4}{\kilo\newton\per\coulomb}$ na direção $y$, produzidos por
duas fontes distintas. Determine o campo resultante.
\resposta{$E=\SI{5}{\kilo\newton\per\coulomb}$, a $\ang{53.1}$ do eixo $x$}
\resolucao{Sendo perpendiculares, compõem-se por Pitágoras:
\[
E=\sqrt{3^{2}+4^{2}}=\SI{5}{\kilo\newton\per\coulomb};
\qquad
\theta=\arctan\frac{4}{3}=\ang{53.1}.
\]
\textbf{Contraste com o potencial:} se esses mesmos pontos tivessem potenciais
de $3$ e $\SI{4}{\kilo\volt}$, o total seria simplesmente
$\SI{7}{\kilo\volt}$ --- soma algébrica, sem geometria.\\
É essa diferença que torna o potencial a ferramenta preferida para cálculos, e o
campo a ferramenta preferida para entender forças e movimento.}
\end{questao}

\begin{questao}[2]{}
Em uma representação por linhas de campo, o que indica uma região onde as
linhas estão mais próximas entre si?
\begin{alternativas}[2]
\item potencial maior
\item campo mais intenso
\item carga maior naquele ponto
\item ausência de campo
\end{alternativas}
\resposta{(b)}
\resolucao{A \textbf{densidade de linhas} é proporcional ao módulo do campo ---
essa é a convenção que dá utilidade quantitativa ao desenho.\\
\textbf{Três regras das linhas de campo:}
\begin{enumerate}
\item Nascem em cargas positivas e morrem em negativas (ou no infinito).
\item São tangentes a $\vec E$ em cada ponto.
\item \textbf{Nunca se cruzam} --- em um cruzamento, o campo teria duas
direções, o que é impossível.
\end{enumerate}
\textbf{Cuidado com (a):} potencial e campo são grandezas distintas. Linhas
próximas indicam campo intenso, o que corresponde a equipotenciais próximas ---
mas não diz nada sobre o \emph{valor} de $V$.}
\end{questao}

\begin{questao}[2]{}
Uma carga produz campo de $\SI{8}{\kilo\newton\per\coulomb}$ no vácuo. O
conjunto é imerso em um dielétrico de $\varepsilon_r=4$. Determine o novo campo.
\resposta{$E=\SI{2}{\kilo\newton\per\coulomb}$}
\resolucao{
\[
E_{meio}=\frac{E_{vácuo}}{\varepsilon_r}=\frac{8}{4}
=\SI{2}{\kilo\newton\per\coulomb}.
\]
\textbf{Mecanismo:} o dielétrico se \emph{polariza} --- suas moléculas se
orientam e criam um campo interno oposto, que cancela parcialmente o campo
original.\\
\textbf{Consequência tecnológica:} o dielétrico reduz o campo para a mesma
carga, o que permite armazenar \emph{mais} carga para a mesma tensão --- é
exatamente por isso que $C=\varepsilon_r C_0$ em capacitores.}
\end{questao}

\begin{questao}[2]{}
Uma gota de massa $\SI{2}{\milli\gram}$ e carga $\SI{1}{\nano\coulomb}$ deve
ficar suspensa em repouso. Determine o campo elétrico necessário.
\dados{$g=\SI{9.8}{\meter\per\second\squared}$}
\resposta{$E=\SI{19.6}{\kilo\newton\per\coulomb}$}
\resolucao{Equilíbrio entre peso e força elétrica:
\[
qE=mg
\;\Longrightarrow\;
E=\frac{mg}{q}=\frac{(\pot{2}{-6})(9{,}8)}{\pot{1}{-9}}
=\pot{1.96}{4}\,\si{\newton\per\coulomb}.
\]
\textbf{Orientação:} o campo deve produzir força \emph{para cima}. Se a gota for
positiva, $\vec E$ aponta para cima; se negativa, para baixo.\\
\textbf{Este é o princípio do experimento de Millikan}, que determinou a carga
elementar. Medindo $E$ e $m$ de gotas suspensas, ele encontrou que todas as
cargas eram múltiplas inteiras de $\pot{1.6}{-19}\,\si{\coulomb}$ --- a
primeira evidência direta da quantização.}
\end{questao}

\blocoaprofundamento

\begin{questao}[3]{}
Duas cargas iguais $Q=+\SI{1}{\micro\coulomb}$ ocupam
$(\pm\SI{0.20}{\meter},\,0)$.
\begin{itensq}
\item Determine o campo em $(0,\ \SI{0.20}{\meter})$.
\item Determine o campo no ponto médio e justifique.
\end{itensq}
\resposta{(a) $E=\SI{159}{\kilo\newton\per\coulomb}$ ao longo de $+y$;
(b) $\vec E=\vec0$}
\resolucao{(a) Cada carga dista
$r=\sqrt{2}\,(0{,}20)=\SI{0.283}{\meter}$:
\[
E_1=E_2=\frac{\pot{9}{9}(\pot{1}{-6})}{0{,}08}
=\pot{1.125}{5}\,\si{\newton\per\coulomb}.
\]
As componentes em $x$ cancelam; as em $y$ somam-se, cada uma multiplicada por
$\cos\ang{45}=0{,}707$:
\[
E=2(\pot{1.125}{5})(0{,}707)=\pot{1.59}{5}\,\si{\newton\per\coulomb}.
\]
(b) \textbf{No ponto médio}, as duas contribuições têm módulos iguais e sentidos
\emph{opostos} (cada carga empurra para longe de si):
\[
\vec E=\vec 0.
\]
\textbf{Mas o potencial ali não é nulo:}
$V=2\ke Q/0{,}20=\SI{90}{\kilo\volt}$. É o caso complementar ao do dipolo, em
que $V=0$ e $\vec E\ne\vec0$ --- e mostra novamente que as duas condições são
independentes.\\
\textbf{Estabilidade:} uma carga positiva no ponto médio está em equilíbrio, mas
\emph{instável} na direção perpendicular ao eixo --- conforme o teorema de
Earnshaw.}
\end{questao}

\begin{questao}[3]{}
Um dipolo é formado por $+q$ e $-q$ separadas por $2a$. Considere um ponto do
eixo perpendicular à linha que as une, a distância $y$ do centro.
\begin{itensq}
\item Deduza a expressão exata do campo.
\item Obtenha a forma assintótica para $y\gg a$.
\item Compare a queda com a de uma carga isolada.
\end{itensq}
\resposta{$E=\dfrac{2\ke qa}{(y^{2}+a^{2})^{3/2}}\to\dfrac{2\ke qa}{y^{3}}$}
\resolucao{(a) Cada carga dista $r=\sqrt{y^{2}+a^{2}}$ e produz
$\ke q/r^{2}$. Por simetria, as componentes ao longo de $y$ se cancelam (uma
aponta para cima, a outra para baixo) e restam as componentes \emph{paralelas ao
eixo do dipolo}, cada uma multiplicada por $a/r$:
\[
E=2\cdot\frac{\ke q}{y^{2}+a^{2}}\cdot\frac{a}{\sqrt{y^{2}+a^{2}}}
=\frac{2\ke qa}{(y^{2}+a^{2})^{3/2}}.
\]
(b) Para $y\gg a$, $(y^{2}+a^{2})^{3/2}\to y^{3}$:
\[
E\approx\frac{2\ke qa}{y^{3}}=\frac{\ke p}{y^{3}},
\]
com $p=2qa$ o momento dipolar.\\
(c) \textbf{O campo do dipolo cai com $1/y^{3}$}, muito mais rápido que o
$1/r^{2}$ de uma carga isolada.\\
\textbf{Razão física:} de longe, as duas cargas quase se cancelam --- o sistema
é \emph{globalmente neutro}. O que sobra é o efeito de sua \emph{separação},
uma correção de ordem superior.\\
\textbf{Padrão geral:} carga isolada $\sim1/r^{2}$; dipolo
$\sim1/r^{3}$; quadrupolo $\sim1/r^{4}$. Quanto mais neutra a distribuição, mais
rápido o campo se extingue --- e é por isso que a matéria comum, neutra, não
exerce força elétrica a distância.}
\end{questao}

\begin{questao}[3]{}
Duas placas condutoras extensas e paralelas têm densidades superficiais
$+\sigma$ e $-\sigma$, com
$\sigma=\SI{1}{\micro\coulomb\per\meter\squared}$.
\begin{itensq}
\item Determine o campo entre as placas.
\item Determine o campo fora delas.
\item Justifique fisicamente.
\end{itensq}
\resposta{Entre: $E=\SI{113}{\kilo\newton\per\coulomb}$; fora: $E=0$}
\resolucao{(a) Cada plano contribui com $\sigma/2\varepsilon_0$, e
\textbf{entre} as placas as duas contribuições têm o \emph{mesmo} sentido (do
positivo para o negativo):
\[
E=2\cdot\frac{\sigma}{2\varepsilon_0}=\frac{\sigma}{\varepsilon_0}
=\frac{\pot{1}{-6}}{\pot{8.85}{-12}}=\pot{1.13}{5}\,\si{\newton\per\coulomb}.
\]
(b) \textbf{Fora}, as duas contribuições têm sentidos \emph{opostos} e módulos
iguais:
\[
E=\frac{\sigma}{2\varepsilon_0}-\frac{\sigma}{2\varepsilon_0}=0.
\]
(c) \textbf{Justificativa:} o campo de um plano infinito é uniforme e aponta
para longe dele (se positivo) ou para ele (se negativo). Na região externa,
qualquer ponto está do mesmo lado das duas placas --- e os campos se opõem.
Entre elas, o ponto está "depois" da positiva e "antes" da negativa, e os campos
se somam.\\
\textbf{Consequência: blindagem natural.} O par de placas confina completamente
seu campo na região interna --- não há campo escapando para fora. É por isso que
um capacitor de placas paralelas não interfere no circuito vizinho, e por que a
energia armazenada está inteiramente entre as armaduras.}
\end{questao}

\begin{questao}[3]{}
Uma esfera condutora de raio $R=\SI{10}{\centi\meter}$ tem carga
$Q=\SI{5}{\micro\coulomb}$.
\begin{itensq}
\item Determine o campo no interior, na superfície e a $\SI{30}{\centi\meter}$
      do centro.
\item Avalie se a configuração é fisicamente sustentável no ar.
\end{itensq}
\resposta{$0$; $\SI{4.5}{\mega\volt\per\meter}$;
$\SI{500}{\kilo\volt\per\meter}$ --- \textbf{insustentável}}
\resolucao{(a) \textbf{Interior:} $E=0$, condição de equilíbrio de qualquer
condutor.\\
\textbf{Superfície:} toda a carga se comporta como concentrada no centro, para
pontos externos:
\[
E=\frac{\ke Q}{R^{2}}=\frac{\pot{9}{9}(\pot{5}{-6})}{0{,}01}
=\pot{4.5}{6}\,\si{\volt\per\meter}.
\]
\textbf{A $\SI{30}{\centi\meter}$:}
$E=\dfrac{\pot{4.5}{4}}{0{,}09}=\pot{5}{5}\,\si{\volt\per\meter}$.\\
(b) \textbf{Insustentável.} O campo na superfície,
$\SI{4.5}{\mega\volt\per\meter}$, excede a rigidez dielétrica do ar
($\SI{3}{\mega\volt\per\meter}$) em $50\%$. O ar ioniza e a carga escapa por
descarga corona.\\
\textbf{Carga máxima real:}
$Q_{\max}=E_{rup}R^{2}/\ke=\SI{3.33}{\micro\coulomb}$. A esfera simplesmente
\emph{não consegue} reter os $\SI{5}{\micro\coulomb}$ enunciados.\\
\textbf{Lição de método:} sempre verifique se o resultado é fisicamente
possível. Um problema pode ser aritmeticamente consistente e descrever uma
situação inexistente.}
\end{questao}

\begin{questao}[3]{}
Um elétron parte do repouso em um campo uniforme de
$E=\SI{1000}{\volt\per\meter}$ e percorre $\SI{2}{\centi\meter}$.
\begin{itensq}
\item Determine a aceleração e o tempo de percurso.
\item Determine a velocidade final e a energia adquirida.
\end{itensq}
\dados{$e=\pot{1.6}{-19}\,\si{\coulomb}$; $m_e=\pot{9.11}{-31}\,\si{\kilogram}$}
\resposta{$a=\pot{1.76}{14}\,\si{\meter\per\second\squared}$;
$t=\SI{15.1}{\nano\second}$; $v=\pot{2.65}{6}\,\si{\meter\per\second}$;
$\SI{20}{\electronvolt}$}
\resolucao{(a)
\[
a=\frac{eE}{m_e}=\frac{(\pot{1.6}{-19})(1000)}{\pot{9.11}{-31}}
=\pot{1.756}{14}\,\si{\meter\per\second\squared};
\]
\[
t=\sqrt{\frac{2d}{a}}=\sqrt{\frac{2(0{,}02)}{\pot{1.756}{14}}}
=\pot{1.51}{-8}\,\si{\second}=\SI{15.1}{\nano\second}.
\]
(b) $v=at=(\pot{1.756}{14})(\pot{1.51}{-8})
=\pot{2.65}{6}\,\si{\meter\per\second}$.\\
\textbf{Energia:} o mais rápido é usar a tensão percorrida,
$U=Ed=1000(0{,}02)=\SI{20}{\volt}$:
\[
W=eU=\SI{20}{\electronvolt}=\pot{3.2}{-18}\,\si{\joule}.
\]
\emph{(Conferindo: $\frac12m_ev^{2}=\frac12(\pot{9.11}{-31})(\pot{2.65}{6})^{2}
=\pot{3.2}{-18}\,\si{\joule}$ \checkmark.)}\\
\textbf{Verificação relativística:} $v/c=0{,}9\%$ --- plenamente no regime
clássico.}
\end{questao}

\begin{questao}[3]{}
Compare a queda do campo com a distância para três configurações: carga
puntiforme, fio infinito carregado e plano infinito carregado.
\begin{itensq}
\item Escreva a dependência de cada uma.
\item Explique a origem geométrica das diferenças.
\end{itensq}
\resposta{$1/r^{2}$, $1/r$ e constante}
\resolucao{(a)
\begin{center}
\begin{tabular}{lcc}
\hline
Distribuição & Campo & Queda\\
\hline
Carga puntiforme & $\dfrac{\ke Q}{r^{2}}$ & $1/r^{2}$\\[6pt]
Fio infinito ($\lambda$) & $\dfrac{\lambda}{2\pi\varepsilon_0 r}$ & $1/r$\\[6pt]
Plano infinito ($\sigma$) & $\dfrac{\sigma}{2\varepsilon_0}$ & constante\\
\hline
\end{tabular}
\end{center}
(b) \textbf{Origem geométrica.} Pela lei de Gauss, o fluxo total é fixo pela
carga envolvida, e o campo é o fluxo dividido pela área da superfície:
\begin{itemize}
\item \textbf{Ponto:} a superfície é uma \emph{esfera}, de área
$4\pi r^{2}$ --- o campo cai com $1/r^{2}$.
\item \textbf{Fio:} a superfície é um \emph{cilindro}, de área lateral
$2\pi r\ell$ --- cresce apenas com $r$, e o campo cai com $1/r$.
\item \textbf{Plano:} a superfície é um \emph{prisma} cuja área \emph{não
depende} de $r$ --- e o campo não cai.
\end{itemize}
\textbf{Regra unificadora:} o campo cai como $1/r^{n}$, onde $n$ é o número de
dimensões em que a fonte é \emph{limitada}. Ponto: limitado em três
($n=2$). Fio: em duas ($n=1$). Plano: em uma ($n=0$).\\
\textbf{Consequência prática:} perto de uma placa carregada extensa, afastar-se
não ajuda --- o campo permanece. É o oposto da intuição formada com cargas
puntiformes.}
\end{questao}

\begin{questao}[3]{}
Uma haste isolante fina é carregada de modo que produz, em certo ponto, campo
de $\SI{600}{\newton\per\coulomb}$. A haste é então partida ao meio, e apenas
uma das metades permanece.
\begin{itensq}
\item O campo cai à metade? Justifique.
\item Determine em que condição isso valeria.
\end{itensq}
\resposta{Não necessariamente --- só vale se a metade removida contribuísse com
exatamente metade do campo}
\resolucao{(a) \textbf{Não, em geral.} O campo é a soma \emph{vetorial} das
contribuições de todos os elementos da haste, e cada elemento contribui com
módulo e direção diferentes, dependendo de sua distância e posição angular em
relação ao ponto.\\
Se o ponto estiver \emph{sobre a mediatriz} da haste, as duas metades
contribuem igualmente e o campo de fato cai à metade --- mas a \emph{direção}
resultante também muda, pois as componentes que antes se cancelavam agora não
se cancelam mais.\\
Se o ponto estiver \emph{sobre o prolongamento} da haste, a metade mais próxima
contribui muito mais que a distante, e remover a distante reduz o campo bem
menos que $50\%$.\\
(b) \textbf{A queda é exatamente à metade quando:}
\begin{itemize}
\item as duas metades contribuem com vetores \emph{idênticos} --- o que exige
simetria completa em relação ao ponto; \emph{e}
\item essa condição só se realiza em pontos muito distantes, onde a haste inteira
se comporta como carga puntiforme e a divisão é simplesmente $Q\to Q/2$.
\end{itemize}
\textbf{Lição:} a superposição vale \emph{sempre}, mas ela é vetorial. Reduzir
a fonte pela metade não reduz o campo pela metade, exceto em casos de simetria
particular.}
\end{questao}

\blocodesafio

\begin{questao}[4]{}
\textbf{(Campo nulo com potencial não nulo.)} Construa uma configuração em que
$\vec E=\vec0$ em um ponto onde $V\ne0$.
\begin{itensq}
\item Proponha o arranjo mais simples e determine os valores.
\item Explique por que isso é possível.
\item Determine se o inverso ($V=0$ com $\vec E\ne\vec0$) também ocorre.
\end{itensq}
\resposta{Duas cargas iguais $+Q$; no ponto médio $\vec E=\vec0$ e
$V=2\ke Q/a$}
\resolucao{(a) \textbf{Arranjo:} duas cargas iguais $Q=+\SI{2}{\micro\coulomb}$
em $(\pm\SI{0.20}{\meter},\,0)$. No ponto médio:
\[
\vec E=\vec 0
\quad\text{(contribuições opostas)};
\]
\[
V=2\cdot\frac{\ke Q}{0{,}20}
=2\cdot\frac{\pot{9}{9}(\pot{2}{-6})}{0{,}20}
=\pot{1.8}{5}\,\si{\volt}=\SI{180}{\kilo\volt}.
\]
Campo nulo, potencial de cento e oitenta mil volts.\\
(b) \textbf{Por que é possível.} As duas grandezas se relacionam por
$\vec E=-\nabla V$: o campo é a \emph{taxa de variação} de $V$, não seu valor.
Uma função pode ter valor grande e derivada nula --- é exatamente o que ocorre
em um ponto \emph{estacionário}.\\
No ponto médio, $V$ tem um mínimo ao longo do eixo que une as cargas: o valor é
alto, mas a inclinação é zero.\\
(c) \textbf{Sim, o inverso também ocorre.} Substituindo uma das cargas por
$-Q$, o ponto médio passa a ter $V=0$ (contribuições opostas se cancelam
algebricamente) e $\vec E\ne\vec0$ (os vetores apontam para o mesmo lado e se
somam).\\
\textbf{Síntese:} $V=0$ e $\vec E=\vec0$ são condições \textbf{logicamente
independentes}. Nenhuma implica a outra, em nenhum dos dois sentidos.}
\end{questao}

\begin{questao}[4]{}
\textbf{(Linhas de campo.)} Analise as linhas de campo de um dipolo elétrico.
\begin{itensq}
\item Descreva seu traçado.
\item Demonstre que linhas de campo nunca se cruzam.
\item Explique o que ocorre em pontos de campo nulo.
\end{itensq}
\resposta{Saem de $+q$ e entram em $-q$; não se cruzam porque $\vec E$ tem
direção única em cada ponto}
\resolucao{(a) \textbf{Traçado.} As linhas saem radialmente de $+q$, curvam-se e
entram em $-q$. A linha que segue o eixo do dipolo é reta; as demais formam
arcos progressivamente mais abertos. Algumas partem para o infinito e retornam.\\
A densidade de linhas é máxima entre as cargas --- onde o campo é mais intenso
--- e decresce com $1/r^{3}$ ao se afastar, conforme demonstrado antes.\\
(b) \textbf{Demonstração.} Suponha, por absurdo, que duas linhas se cruzem em
um ponto $P$. Por definição, a linha de campo é \emph{tangente} a $\vec E$ em
cada ponto. Se duas linhas distintas passam por $P$, elas têm tangentes
distintas ali --- e portanto $\vec E(P)$ teria \textbf{duas direções}
simultaneamente.\\
Mas $\vec E$ é uma função vetorial \emph{bem definida}: a cada ponto do espaço
corresponde \emph{um} vetor. Contradição. Logo, as linhas não se cruzam.\\
\textbf{Corolário:} o mesmo argumento vale para linhas de corrente, linhas de
campo magnético e qualquer outro campo vetorial.\\
(c) \textbf{Pontos de campo nulo são a exceção.} Onde $\vec E=\vec0$, não há
direção definida --- e ali as linhas \emph{podem} se encontrar, sem contradição.
São os \textbf{pontos de sela} do potencial.\\
Em um dipolo ($+q$ e $-q$) não existe tal ponto no espaço finito. Mas em duas
cargas \emph{iguais}, o ponto médio é exatamente um ponto de campo nulo --- e
ali as linhas se aproximam de todos os lados sem se tocarem, definindo uma
separatriz.}
\end{questao}

\begin{questao}[4]{}
\textbf{(Projeto de acelerador de placas.)} Projete um arranjo de placas
paralelas que imprima aceleração de
$a=\pot{2}{12}\,\si{\meter\per\second\squared}$ a um elétron.
\begin{itensq}
\item Determine o campo necessário.
\item Determine a tensão para separação de $\SI{1}{\centi\meter}$.
\item Determine o tempo de trânsito e a velocidade de saída.
\end{itensq}
\dados{$e=\pot{1.6}{-19}\,\si{\coulomb}$; $m_e=\pot{9.11}{-31}\,\si{\kilogram}$}
\resposta{$E=\SI{11.4}{\volt\per\meter}$; $U=\SI{0.114}{\volt}$;
$t=\SI{100}{\nano\second}$; $v=\pot{2}{5}\,\si{\meter\per\second}$}
\resolucao{(a)
\[
E=\frac{m_ea}{e}=\frac{(\pot{9.11}{-31})(\pot{2}{12})}{\pot{1.6}{-19}}
=\SI{11.4}{\volt\per\meter}.
\]
(b) $U=Ed=11{,}4(0{,}01)=\SI{0.114}{\volt}$.\\
(c)
\[
t=\sqrt{\frac{2d}{a}}=\sqrt{\frac{0{,}02}{\pot{2}{12}}}
=\pot{1}{-7}\,\si{\second}=\SI{100}{\nano\second};
\]
\[
v=at=(\pot{2}{12})(\pot{1}{-7})=\pot{2}{5}\,\si{\meter\per\second}.
\]
\textbf{O resultado é surpreendente e instrutivo.} Uma aceleração de
$\pot{2}{12}\,\si{\meter\per\second\squared}$ --- duzentos bilhões de vezes
$g$ --- exige apenas $\SI{0.114}{\volt}$, tensão menor que a de uma pilha
descarregada.\\
\textbf{Razão:} a razão carga/massa do elétron é enorme
($\pot{1.76}{11}\,\si{\coulomb\per\kilogram}$). Campos elétricos modestos
produzem acelerações colossais em partículas leves.\\
\textbf{Ressalva prática:} com tensão tão baixa, o projeto seria dominado por
efeitos parasitas --- diferenças de função-trabalho entre os metais das placas
(décimos de volt), cargas superficiais e ruído térmico. Um projeto real usaria
tensões de dezenas de volts e aceleração proporcionalmente maior.}
\end{questao}

\begin{questao}[4]{}
\textbf{(Cavidade em condutor.)} Uma carga puntiforme $+q$ é colocada no
interior de uma cavidade vazia dentro de um condutor \emph{neutro} e isolado.
\begin{itensq}
\item Determine a carga induzida em cada superfície.
\item Determine o campo dentro do metal e fora do condutor.
\item Analise o que muda se o condutor for aterrado.
\end{itensq}
\resposta{$-q$ na parede da cavidade, $+q$ na superfície externa; campo nulo no
metal e não nulo fora}
\resolucao{(a) \textbf{Parede da cavidade:} aplicando a lei de Gauss a uma
superfície inteiramente dentro do metal (onde $\vec E=\vec0$), o fluxo é nulo e
a carga envolvida também. Como a superfície envolve $+q$ e a parede interna:
\[
q+q_{interna}=0
\;\Longrightarrow\;
q_{interna}=-q.
\]
\textbf{Superfície externa:} o condutor era neutro, então a carga total deve
permanecer nula:
\[
q_{externa}=+q.
\]
(b) \textbf{Dentro do metal:} $\vec E=\vec0$, por equilíbrio.\\
\textbf{Fora:} a carga $+q$ da superfície externa produz campo --- e ele é
\textbf{esfericamente simétrico} se a superfície externa for esférica,
\emph{independentemente} de onde a carga esteja dentro da cavidade.\\
\textbf{Consequência notável:} de fora, é impossível saber a posição da carga
interna, nem a forma da cavidade. O condutor "apaga" essa informação.\\
(c) \textbf{Se o condutor for aterrado:} a carga $+q$ da superfície externa
escoa para a terra. Resta apenas $-q$ na parede da cavidade, e
\[
\vec E=\vec 0\quad\text{em todo o exterior}.
\]
\textbf{A blindagem torna-se completa nos dois sentidos.} É exatamente esse o
princípio da gaiola de Faraday aterrada --- e a diferença entre blindar contra
campos externos (que dispensa aterramento) e impedir que fontes internas
irradiem (que o exige).}
\end{questao}

\begin{questao}[4]{}
\textbf{(Esfera isolante uniformemente carregada.)} Uma esfera isolante maciça
de raio $R$ tem carga $Q$ distribuída uniformemente em seu \emph{volume}.
\begin{itensq}
\item Determine o campo para $r<R$ e para $r>R$.
\item Identifique onde o campo é máximo.
\item Compare com a esfera \emph{condutora} de mesma carga.
\end{itensq}
\resposta{$E=\dfrac{\ke Qr}{R^{3}}$ dentro e $\dfrac{\ke Q}{r^{2}}$ fora;
máximo na superfície}
\resolucao{(a) \textbf{Fora ($r>R$):} pela lei de Gauss com superfície esférica,
toda a carga está envolvida:
\[
E=\frac{\ke Q}{r^{2}},
\]
idêntico ao de uma carga puntiforme no centro.\\
\textbf{Dentro ($r<R$):} apenas a fração de carga interna à superfície gaussiana
conta. Como a densidade é uniforme, essa fração é $(r/R)^{3}$:
\[
E=\frac{\ke\,Q(r/R)^{3}}{r^{2}}=\frac{\ke Q\,r}{R^{3}},
\]
\textbf{proporcional a $r$} --- cresce linearmente do centro à superfície.\\
(b) \textbf{Máximo na superfície} ($r=R$), onde as duas expressões coincidem em
$E=\ke Q/R^{2}$. O campo é contínuo ali, embora com derivadas diferentes de
cada lado.\\
\emph{No centro, $E=0$ por simetria.}\\
(c) \textbf{Comparação com a esfera condutora.}
\begin{center}
\begin{tabular}{lcc}
\hline
Região & Isolante & Condutora\\
\hline
$r<R$ & $\ke Qr/R^{3}$ & $0$\\
$r=R$ & $\ke Q/R^{2}$ & $\ke Q/R^{2}$\\
$r>R$ & $\ke Q/r^{2}$ & $\ke Q/r^{2}$\\
\hline
\end{tabular}
\end{center}
\textbf{São indistinguíveis do lado de fora} --- a lei de Gauss só "enxerga" a
carga total envolvida, não sua distribuição. A diferença está inteiramente no
interior: no condutor a carga migra para a superfície e o campo interno é nulo;
no isolante ela permanece onde foi colocada.}
\end{questao}

\begin{questao}[4]{}
\textbf{(Traçado numérico de linhas de campo.)} Descreva um procedimento
computacional para traçar linhas de campo a partir de uma expressão conhecida de
$\vec E(x,y)$.
\begin{itensq}
\item Formule o algoritmo.
\item Indique como escolher os pontos de partida.
\item Aponte as dificuldades numéricas e como contorná-las.
\end{itensq}
\resposta{Integrar $\dfrac{\mathrm{d}\vec r}{\mathrm{d}s}
=\dfrac{\vec E}{|\vec E|}$ a partir de pontos distribuídos em torno das cargas}
\resolucao{(a) \textbf{Algoritmo.} A linha de campo é, por definição, tangente a
$\vec E$. Parametrizando pelo comprimento de arco $s$:
\[
\frac{\mathrm{d}\vec r}{\mathrm{d}s}=\frac{\vec E(\vec r)}{|\vec E(\vec r)|}.
\]
Numericamente, com passo $h$:
\begin{enumerate}
\item Escolha um ponto inicial $\vec r_0$.
\item Calcule $\vec E(\vec r_n)$ e normalize.
\item Avance: $\vec r_{n+1}=\vec r_n+h\,\hat E$ (Euler) ou, melhor, por
Runge--Kutta de quarta ordem.
\item Repita até sair do domínio ou atingir uma carga negativa.
\end{enumerate}
\textbf{Normalizar é essencial:} sem isso, o passo seria proporcional a
$|\vec E|$ e a linha "voaria" perto das cargas e travaria longe delas.\\
(b) \textbf{Pontos de partida.} Distribua-os em pequenos círculos ao redor de
cada carga \emph{positiva}, com espaçamento angular uniforme. Para que a
densidade de linhas represente corretamente a intensidade, o número de linhas
por carga deve ser \emph{proporcional a $|q|$} --- oito linhas para $+q$,
dezesseis para $+2q$.\\
(c) \textbf{Dificuldades e soluções.}
\begin{itemize}
\item \textbf{Singularidades:} $|\vec E|\to\infty$ nas cargas. Solução: iniciar
a uma distância pequena mas finita, e interromper a integração ao entrar em um
raio mínimo em torno de uma carga negativa.
\item \textbf{Pontos de campo nulo:} ali $\hat E$ é indefinido e o algoritmo
trava ou oscila. Solução: detectar $|\vec E|<\epsilon$ e encerrar a linha.
\item \textbf{Acúmulo de erro:} o método de Euler desvia sistematicamente em
regiões de curvatura alta. Solução: Runge--Kutta e passo adaptativo, menor onde
o campo varia rápido.
\item \textbf{Linhas que escapam:} defina um domínio e encerre ao ultrapassá-lo.
\end{itemize}
\textbf{Verificação do resultado:} as linhas traçadas devem ser perpendiculares
às equipotenciais calculadas independentemente --- teste que detecta quase todo
erro de implementação.}
\end{questao}

\begin{questao}[4]{}
\textbf{(Limite de campo em capacitor.)} Um capacitor de placas paralelas com
$d=\SI{0.1}{\milli\meter}$ usa dielétrico de rigidez
$\SI{20}{\kilo\volt\per\milli\meter}$ e $\varepsilon_r=3$.
\begin{itensq}
\item Determine a tensão máxima e o campo correspondente.
\item Determine a densidade superficial de carga no limite.
\item Discuta o compromisso entre capacitância e tensão.
\end{itensq}
\dados{$\varepsilon_0=\pot{8.85}{-12}\,\si{\farad\per\meter}$}
\resposta{$U_{\max}=\SI{2}{\kilo\volt}$;
$E=\SI{2e7}{\volt\per\meter}$;
$\sigma=\SI{531}{\micro\coulomb\per\meter\squared}$}
\resolucao{(a)
\[
U_{\max}=E_{rup}\,d=20\times0{,}1=\SI{2}{\kilo\volt};
\qquad
E=\pot{2}{7}\,\si{\volt\per\meter}.
\]
(b) No dielétrico, $E=\dfrac{\sigma}{\varepsilon_r\varepsilon_0}$, logo
\[
\sigma=\varepsilon_r\varepsilon_0E
=3(\pot{8.85}{-12})(\pot{2}{7})
=\pot{5.31}{-4}\,\si{\coulomb\per\meter\squared}.
\]
(c) \textbf{O compromisso.} A capacitância por unidade de área é
$C/A=\varepsilon_r\varepsilon_0/d$: reduzir $d$ \emph{aumenta} $C$. Mas
$U_{\max}=E_{rup}d$: reduzir $d$ \emph{diminui} a tensão suportável.\\
O produto revela o que realmente importa:
\[
\frac{C\,U_{\max}}{A}=\varepsilon_r\varepsilon_0E_{rup},
\]
que \textbf{não depende de $d$}. E a energia por unidade de volume,
\[
u=\frac12\varepsilon_r\varepsilon_0E_{rup}^{2},
\]
também não. \textbf{A geometria não decide nada} --- só o material.\\
\textbf{Consequência:} não existe truque de projeto que aumente a energia
armazenada por volume. Todo avanço em capacitores vem de materiais com
$\varepsilon_r$ ou $E_{rup}$ maiores --- e esses dois parâmetros costumam ser
antagônicos, o que torna o problema difícil.}
\end{questao}

\begin{questao}[4]{}
\textbf{(Superposição vetorial em anel.)} Um anel de raio $R$ tem carga
$Q$ uniformemente distribuída.
\begin{itensq}
\item Explique por que o campo é nulo no centro.
\item Determine, sem integrar, o campo em um ponto muito distante no eixo.
\item Discuta o que ocorre se metade do anel for removida.
\end{itensq}
\resposta{Centro: $\vec E=\vec0$ por simetria; longe: $E\to\ke Q/x^{2}$;
com meio anel, campo não nulo no centro}
\resolucao{(a) \textbf{No centro,} para cada elemento de carga existe outro
diametralmente oposto, à mesma distância $R$. Suas contribuições têm módulos
iguais e sentidos opostos, cancelando-se aos pares:
\[
\vec E_{centro}=\vec 0.
\]
\textbf{Note:} o \emph{potencial} no centro não é nulo --- vale
$V=\ke Q/R$, pois potenciais somam-se sem cancelamento entre contribuições de
mesmo sinal. Mais um caso de $\vec E=\vec0$ com $V\ne0$.\\
(b) \textbf{Muito longe} ($x\gg R$), todos os elementos estão praticamente à
mesma distância $x$ e praticamente na mesma direção. O anel se comporta como uma
carga puntiforme:
\[
E\to\frac{\ke Q}{x^{2}}.
\]
\textbf{Este é o teste que toda expressão deve passar:} qualquer distribuição
finita, vista de longe, deve reproduzir a lei de Coulomb.\\
(c) \textbf{Com meio anel}, a simetria de cancelamento se perde. Cada elemento
não tem mais seu par oposto, e o campo no centro \textbf{deixa de ser nulo}.\\
Pelo argumento da superposição ao contrário: como o anel inteiro dá campo nulo,
o campo do semianel restante é o \emph{negativo} do campo do semianel removido
--- ou seja, os dois têm o mesmo módulo. Integrando (como no caso do fio
semicircular), obtém-se
\[
E=\frac{2\ke Q'}{\pi R^{2}},
\]
com $Q'=Q/2$ a carga de cada metade, dirigido ao longo do eixo de simetria.}
\end{questao}

% END BANCO COMPLEMENTAR

\gabarito[1.3 Campo elétrico]
